% joint.tex -- first full draft of the joint paper (G. Araujo-Pardo, N. Lopez, D. Ozden).
% Assembled from the lab's refereed notes (see the source list in the comment at the end of this file).
% Build: tectonic joint.tex
\documentclass[11pt]{amsart}
\usepackage[margin=2.7cm]{geometry}
\usepackage{amsmath,amssymb,amsthm,mathtools}
\usepackage{array,booktabs,longtable}
\usepackage{enumitem}
\usepackage[dvipsnames]{xcolor}
\usepackage[colorinlistoftodos,textsize=small]{todonotes}
\presetkeys{todonotes}{inline,backgroundcolor=yellow!25,bordercolor=orange}{}
\usepackage[hidelinks]{hyperref}
\usepackage{url}

\setlength{\LTcapwidth}{\textwidth}
\makeatletter\providecommand\@dotsep{4.5}\makeatother
\setlist{leftmargin=2.2em}

\newtheorem{theorem}{Theorem}[section]
\newtheorem{lemma}[theorem]{Lemma}
\newtheorem{proposition}[theorem]{Proposition}
\newtheorem{corollary}[theorem]{Corollary}
\theoremstyle{definition}
\newtheorem{definition}[theorem]{Definition}
\newtheorem{construction}[theorem]{Construction}
\newtheorem{example}[theorem]{Example}
\newtheorem{fact}[theorem]{Fact}
\theoremstyle{remark}
\newtheorem{remark}[theorem]{Remark}

\numberwithin{equation}{section}

\newcommand{\F}{\mathbb F}
\newcommand{\Z}{\mathbb Z}
\newcommand{\PG}{\mathrm{PG}}
\newcommand{\AS}{\mathrm{AS}}
\newcommand{\GF}{\mathrm{GF}}
\newcommand{\Tr}{\mathrm{Tr}}
\newcommand{\Nm}{\mathrm{N}}
\newcommand{\I}{\mathrel{\mathrm I}}
\newcommand{\nb}{n_b}
\newcommand{\Mb}{M_b}
\newcommand{\diam}{\operatorname{diam}}
\newcommand{\ecc}{\operatorname{ecc}}
\newcommand{\opp}{\operatorname{opp}}
\newcommand{\mdp}{\operatorname{mid}}
\newcommand{\cA}{\mathcal A}
\newcommand{\cB}{\mathcal B}
\newcommand{\cE}{\mathcal E}
\newcommand{\cL}{\mathcal L}
\newcommand{\cM}{\mathcal M}
\newcommand{\cO}{\mathcal O}
\newcommand{\cP}{\mathcal P}
\newcommand{\cS}{\mathcal S}
\newcommand{\eps}{\varepsilon}
\newcommand{\SL}{\mathrm{SL}}
\newcommand{\PSL}{\mathrm{PSL}}
\newcommand{\Om}{\Omega\mathrm H}

\title[Large bipartite graphs of given degree and diameter]{Large bipartite graphs of given degree and diameter\\
from generalized polygons}

\author{Gabriela Araujo-Pardo}
\author{Nacho L\'opez}
\author{Do\u{g}a \"Ozden}
\date{Draft of 28 October 2026}

\begin{document}

\begin{abstract}
Let $\nb(\Delta,D)$ be the largest order of a bipartite graph of maximum degree at most $\Delta$ and diameter at most
$D$. We give constructions from generalized polygons that improve the lower bounds in the 2013 table of Miller and
\v{S}ir\'a\v{n} at diameters 3, 4, 6 and 7, and Haar graphs that improve it at diameters 5, 8, 9 and 10. At diameter 3
we add a fourth hat to the three hats of Araujo-Pardo and L\'opez and obtain $\nb(q+2,3)\ge 2(q^2+q+6)$ with regular
graphs. At diameter 4 we copy points and lines of a generalized quadrangle with weights and stand-ins; a weighted
2-ovoid of the Ahrens--Szekeres quadrangle built from two elliptic quadrics gives $\nb(q+2,4)\ge 2q^2(q+2)$ for odd
$q$, and Singer 2-ovoids of $W(3,q)$ give $\nb(q+3,4)\ge 2(q^2+1)(q+3)$, both with regular graphs. At diameter 6 an
explicit set of $q^2(q^2-1)$ lines of the split Cayley hexagon $H(q)$ (``whiskers'') added to Delorme's doubling gives
$\nb(q+m,6)\ge 2(q+1)(q^4+q^2+1)+(m-1)\bigl(2q^4+2q^3+2q+2+m(q^3+q+2)+m^2\bigr)$, and a regular version. At diameter 7
we join $2r+1$ copies of the incidence graph of $H(q)$ by matchings, delete matching edges at suitable sets and attach
new vertices; we prove which sets are allowed. All constructions are given in projective coordinates with the exact
neighbourhood of every new vertex. We also show that the graphs of G\'omez behind six diameter-8 entries of the table
have diameter at least 9 as defined, and that relabelling their copies by automorphisms does not help.
\end{abstract}

\maketitle

\makeatletter\@starttoc{tdo}{Points for the co-authors}\makeatother

\todo{Title: a placeholder. Choose the final title.}

\todo{Authors, order, affiliations, e-mail addresses. The byline of [AL20] is ``Gabriela Araujo, Nacho L\'opez''; use
whichever form you prefer here and in the references.}

\todo{AI disclosure. The constructions, proofs and checks in this draft were produced in a research repository run by
Claude Code agents coordinated by D.\ \"Ozden, and most explicit ingredients (sets, connection sets, hub graphs) were
found by computer search. Following the guidance of the 20 October meeting this draft keeps search and AI out of the
mathematics and mentions computation only as ``computer-checked''. Decide how to acknowledge the AI use (most journals
now ask for a statement) and whether to cite the repository \url{https://github.com/dogaozden/degree-diameter}, which
holds every graph, builder and check log.}

%======================================================================
\section{Introduction}\label{sec:intro}
%======================================================================

A bipartite graph of maximum degree $\Delta\ge2$ and diameter $D\ge2$ has at most
\[
\Mb(\Delta,D)=2\bigl(1+(\Delta-1)+\dots+(\Delta-1)^{D-1}\bigr)
\]
vertices (the bipartite Moore bound). Let $\nb(\Delta,D)$ be the largest order of such a graph. For $D\ge3$ and
$\Delta\ge3$ the bound is attained only by the incidence graphs of generalized polygons: projective planes ($D=3$),
generalized quadrangles ($D=4$) and generalized hexagons ($D=6$) of order $(\Delta-1,\Delta-1)$
\cite{FeitHigman,MS13}. These are known when $\Delta-1$ is a prime power. For the other cells only lower bounds are
known. They are collected in Table~4 of the survey of Miller and \v{S}ir\'a\v{n} \cite{MS13} (version of 2013), which
the CombinatoricsWiki \cite{CW} reproduces; since 2013 the wiki has changed only at $(13,3)$ and $(15,3)$, by the work of
Araujo-Pardo and L\'opez \cite{AL20}. We call this table ``the 2013 table''.

This paper gives constructions that improve the 2013 table in the cells listed in Table~\ref{tab:summary}. Most of them
are built from generalized polygons in their classical coordinates, by three kinds of operations: copying points and
lines (with Delorme's doubling with levels \cite{D85a} as the basic tool), changing a bounded number of incidences, and
joining several copies of an incidence graph by matchings and new vertices. The remaining cells are Haar graphs
(bi-Cayley graphs) of explicit groups.

\subsection*{Results}

\begin{itemize}
\item \emph{Diameter 3} (Section~\ref{sec:d3}). Araujo-Pardo and L\'opez \cite{AL20} added three ``hats'' to
  $\PG(2,q)$ and obtained $\nb(q+2,3)\ge2(q^2+q+5)$. We add a fourth hat (Theorem~\ref{thm:hat4}):
  $\nb(q+2,3)\ge2(q^2+q+6)$ for every prime power $q\ge4$, attained by a $(q+2)$-regular graph; the construction works
  in every projective plane of order $q\ge3$.
\item \emph{Diameter 4} (Section~\ref{sec:d4}). Theorem~\ref{thm:B}: in any generalized quadrangle, clones of points
  with arbitrary weights and copies of lines with ``stand-ins'' keep diameter 4. With a weighted 2-ovoid of the
  Ahrens--Szekeres quadrangle $\AS(q)$ made of two elliptic quadrics (Theorem~\ref{thm:A}) this gives
  $\nb(q+2,4)\ge2q^2(q+2)$ for odd $q$; with the Singer 2-ovoids of $W(3,q)$ (Proposition~\ref{prop:C}) it gives
  $\nb(q+3,4)\ge2(q^2+1)(q+3)$ for odd $q$; both with regular graphs. For every prime power $q\le\Delta-1$,
  $\nb(\Delta,4)\ge2q^2\Delta$ (Corollary~\ref{cor:B2}).
\item \emph{Diameter 6} (Section~\ref{sec:d6}). A set of $q^2(q^2-1)$ lines of the split Cayley hexagon $H(q)$, given by
  an explicit rule (Theorem~\ref{thm:whiskers}), can be added at level 0 to Delorme's doubling of $H(q)$ along a
  Hermitian spread. This gives Theorem~\ref{thm:D}, and deleting repeated lines and completing gives regular graphs
  (Theorem~\ref{thm:R}).
\item \emph{Diameter 7} (Section~\ref{sec:d7}). The 2013 values at $(7,7)$ and $(9,7)$--$(16,7)$ are the orders
  $(2r+1)\cdot2|H(q)|$ of $2r+1$ copies of the incidence graph of $H(q)$ joined by matchings. We delete matching
  edges at suitable sets of elements and attach new vertices to the freed slots. Theorem~\ref{thm:assembly} reduces the
  diameter to one condition per pair of copies; Theorems~\ref{thm:autotrans} and~\ref{thm:localdual} give two
  sufficient conditions (segment sets at automorphic transitions; local dualities), the first also necessary. With
  classical sets this improves all nine cells (Corollaries~\ref{cor:d7-91}--\ref{cor:d7-r}); explicit larger sets give
  the values of Table~\ref{tab:summary}.
\item \emph{Diameter 8} (Section~\ref{sec:d8}). Six entries of the 2013 table at $D=8$ are the orders of the compound
  graphs $\Om_q$ of G\'omez \cite{Go86}, for which the thesis states diameter at most 8 without proof. As defined there,
  these graphs have diameter at least 9 (Theorems~\ref{thm:proj} and~\ref{thm:twolabels}), and relabelling the copies
  by automorphisms of the hexagon does not give diameter 8 (Theorems~\ref{thm:geomq8} and~\ref{thm:geomq7}).
\item \emph{Haar graphs} (Section~\ref{sec:haar}). We state the covering criterion that decides the diameter of a Haar
  graph from two breadth-first searches, and list Haar graphs that improve two cells at $D=5$ and 34 cells at
  $D=8,9,10$. These are single explicit graphs, checked by computer.
\end{itemize}

Problem~4 of \cite{AL20} asks to generalize the hat construction to regular bipartite graphs of diameters 4 and 6. The
hat lemma (Lemma~\ref{lem:hat}) says in what sense it generalizes, and Corollaries~\ref{cor:A1} and~\ref{cor:C1} and
Theorem~\ref{thm:R} give regular graphs at $D=4$ and $D=6$ that beat the 2013 table.

\subsection*{What is proved and what is computed}
Every theorem in Sections~\ref{sec:prelim}--\ref{sec:d8} has a complete proof here, with one kind of exception, always
marked: a statement about one explicit finite object (a set of points, a connection set, a hub graph, a count of
route words) is \emph{computer-checked}, and we name the check. ``Computer-checked'' for a graph means: built from its
definition, and found to be bipartite, of the stated maximum degree and of the stated diameter by breadth-first search
from every vertex (or from one vertex of every orbit of an automorphism group that was itself checked), by at least two
programs written separately. Appendix~\ref{app:checks} lists the checks. Many of the statements are also proved in
Lean~4 with Mathlib; Appendix~\ref{app:lean} says exactly which.

\begin{longtable}{@{}l r r r l l@{}}
\caption{The cells improved in this paper. $\Mb$: bipartite Moore bound; 2013: \cite[Table 4]{MS13}; ``reg.'': the
largest regular graph we give, when the record graph is not regular. Values at $D=8,9,10$ are in
Table~\ref{tab:haar}.}\label{tab:summary}\\
\toprule
$(\Delta,D)$ & $\Mb$ & 2013 & this paper & reg. & construction \\
\midrule\endfirsthead
\toprule
$(\Delta,D)$ & $\Mb$ & 2013 & this paper & reg. & construction \\
\midrule\endhead
\bottomrule\endfoot
(11,3) & 222 & 190 & \textbf{192} & & fourth hat, $q=9$ (Thm.~\ref{thm:hat4}) \\
(13,3) & 314 & 270\,$^a$ & \textbf{276} & & fourth hat, $q=11$ \\
(15,3) & 422 & 370\,$^a$ & \textbf{376} & & fourth hat, $q=13$ \\
(7,4) & 518 & 346 & \textbf{359} & 350 & clones and shadows on $Q(4,5)$ (Ex.~\ref{ex:359}) \\
(11,4) & 2\,222 & 1\,734 & \textbf{1\,782} & & $\AS(9)$, two quadrics (Cor.~\ref{cor:A1}) \\
(13,4) & 3\,770 & 3\,064 & \textbf{3\,146} & & $\AS(11)$ \\
(15,4) & 5\,910 & 4\,946 & \textbf{5\,070} & & $\AS(13)$ \\
(16,4) & 7\,232 & 5\,134 & \textbf{5\,440} & & Singer 2-ovoid of $W(3,13)$ (Cor.~\ref{cor:C1}) \\
(7,5) & 3\,110 & 1\,134 & \textbf{1\,152} & & Haar graph (Sec.~\ref{sec:haar}) \\
(9,5) & 9\,362 & 2\,496 & \textbf{2\,662} & & Haar graph \\
(7,6) & 18\,662 & 8\,992 & \textbf{9\,592} & 9\,348 & whiskers in $H(5)$ (Thm.~\ref{thm:D}, \ref{thm:R}) \\
(11,6) & 222\,222 & 142\,464 & \textbf{148\,944} & 147\,500 & $H(9)$ \\
(13,6) & 542\,906 & 374\,452 & \textbf{388\,972} & 386\,328 & $H(11)$ \\
(15,6) & 1\,158\,390 & 842\,048 & \textbf{870\,440} & 866\,068 & $H(13)$ \\
(16,6) & 1\,627\,232 & 884\,062 & \textbf{940\,846} & 927\,724 & $H(13)$, $m=3$ \\
(7,7) & 111\,974 & 23\,436 & \textbf{23\,598} & & 3 levels of $H(5)$, local duality (Sec.~\ref{sec:d7}) \\
(9,7) & 599\,186 & 117\,648 & \textbf{118\,238} & & 3 levels of $H(7)$ \\
(10,7) & 1\,195\,742 & 224\,694 & \textbf{225\,152} & & 3 levels of $H(8)$ \\
(11,7) & 2\,222\,222 & 398\,580 & \textbf{413\,180} & & 3 levels of $H(9)$, duality \\
(12,7) & 3\,897\,434 & 664\,300 & \textbf{693\,500} & & 5 levels of $H(9)$, duality \\
(13,7) & 6\,514\,874 & 1\,062\,936 & \textbf{1\,064\,472} & & 3 levels of $H(11)$ \\
(14,7) & 10\,458\,086 & 1\,771\,560 & \textbf{1\,789\,588} & & 5 typed levels of $H(11)$ \\
(15,7) & 16\,217\,462 & 2\,480\,184 & \textbf{2\,501\,491} & & 7 typed levels of $H(11)$ \\
(16,7) & 24\,408\,482 & 4\,022\,340 & \textbf{4\,052\,108} & & 5 typed levels of $H(13)$ \\
\end{longtable}
\noindent $^a$ The current values of \cite{CW} are 274 and 376 \cite{AL20}.

All graphs in Table~\ref{tab:summary} except the four Haar graphs at $D=5$ and $D=8$--$10$ are members of families
proved here for every prime power $q$ (or every odd $q$), with explicit ingredients; at $D=4$ (7,4) and at $D=7$ the
values in the table use, in addition, explicit sets that are computer-checked, and Section~\ref{sec:d7} also gives
values beating the 2013 table in all nine cells from classical ingredients only (Table~\ref{tab:d7general}).

\todo{Scope. The repository also holds graphs at other cells that are not covered by the refereed notes this draft is
built from: $(16,3)$ 396; $(4,5)$ 162, $(5,5)$ 354, $(8,5)$ 2\,048, $(9,5)$ 2\,758, $(10,5)$ 4\,800, $(11,5)$ 7\,020,
$(12,5)$ 9\,840, $(13,5)$ 11\,700, $(14,5)$ 17\,568, $(15,5)$ 20\,736, $(16,5)$ 33\,792 (lifts of generalized
quadrangles and other lifts); $(4,7)$ 864, $(5,7)$ 3\,120, $(6,7)$ 8\,880, $(8,7)$ 43\,344; $(3,10)$ 810; and later
Haar graphs at $(9,8)$ 546\,720, $(5,9)$ 29\,756, $(7,10)$ 2\,298\,240, $(8,10)$ 8\,043\,840. Decide whether this paper
should include them (they would need write-ups and referee reports first), mention them in a remark, or leave them for
a second paper.}

\subsection*{Organisation}
Section~\ref{sec:prelim} fixes notation, coordinates for the classical polygons, and Delorme's doubling with levels,
with a proof and the hat lemma. Sections~\ref{sec:d3}--\ref{sec:d7} treat diameters 3, 4, 6 and 7, Section~\ref{sec:d8}
the G\'omez entries at diameter 8, and Section~\ref{sec:haar} Haar graphs. Section~\ref{sec:table} collects the
updated values. Appendix~\ref{app:359} writes out the $(7,4)$ example, Appendix~\ref{app:lean} lists what is checked in
Lean, and Appendix~\ref{app:checks} the computer checks.

%======================================================================
\section{Preliminaries}\label{sec:prelim}
%======================================================================

\subsection{Incidence structures}
An \emph{incidence structure} $\cS=(\cP,\cL,\I)$ has points, lines and a symmetric incidence relation; repeated points
and repeated lines (with the same incidences) are allowed. Its \emph{incidence graph} $\Gamma(\cS)$ has vertex set
$\cP\cup\cL$ and the incident pairs (\emph{flags}) as edges; it is bipartite with the points on one side. A non-incident
point--line pair is an \emph{anti-flag}. Every bipartite graph with a chosen 2-colouring is of this form, so
\[
\nb(\Delta,D)=\max\{|\cP|+|\cL|:\ \text{every element has at most }\Delta\text{ neighbours},\ \diam\Gamma(\cS)\le D\}.
\]
Distances $d(x,y)$ are taken in $\Gamma(\cS)$, and $\diam\cS$ means $\diam\Gamma(\cS)$. Two distinct points on a
common line are \emph{collinear}; two distinct lines on a common point are \emph{concurrent}. $\F=\GF(q)$ is the field
of order $q$, and for $q$ odd $\chi$ is its quadratic character, with $\chi(0)=0$.

\begin{lemma}[diameter in geometric terms]\label{lem:diamgeom}
Let $\cS$ have no isolated element and let $D\ge2$.
\begin{enumerate}[label=(\roman*)]
\item If $D$ is even, $\diam\cS\le D$ if and only if $d(p,L)\le D-1$ for every anti-flag $(p,L)$.
\item If $D$ is odd, $\diam\cS\le D$ if and only if any two points, and any two lines, are at distance at most $D-1$.
\end{enumerate}
In particular, $\diam\cS\le3$ if and only if any two points are collinear and any two lines are concurrent, and
$\diam\cS\le4$ if and only if for every anti-flag $(p,L)$ some point of $L$ is collinear with $p$.
\end{lemma}

\begin{proof}
(i) Let $p\ne p'$ be points and $L$ a line on $p'$. If $p\I L$ then $d(p,p')=2$; otherwise
$d(p,p')\le d(p,L)+1\le D$. Lines likewise. (ii) For an anti-flag $(p,L)$ take $p'\I L$; then
$d(p,L)\le d(p,p')+1\le D$. The converses hold by parity.
\end{proof}

\begin{lemma}[adding flags, deleting twins]\label{lem:flags}
(a) Adding flags to an incidence structure keeps its incidence graph bipartite and does not increase any distance.
(b) If two lines have the same set of points, deleting one of them does not change any distance between the remaining
elements.
\end{lemma}

\begin{proof}
(a) is clear. (b) A shortest path through the deleted line can use its twin instead; it does not already contain the
twin, since then it could be shortened.
\end{proof}

\subsection{Generalized polygons}
A \emph{generalized $n$-gon} is an incidence structure whose incidence graph has diameter $n$ and girth $2n$; it has
\emph{order $(s,t)$} if every line has $s+1$ points and every point is on $t+1$ lines. It is \emph{thick} if
$s,t\ge2$. A generalized 3-gon of order $(q,q)$ is a projective plane of order $q$, with $q^2+q+1$ points and as many
lines. A generalized quadrangle (GQ) of order $(s,t)$ has $(s+1)(st+1)$ points and $(t+1)(st+1)$ lines, and its axiom
is: for every anti-flag $(p,L)$ exactly one point of $L$ is collinear with $p$. A generalized hexagon of order $(q,q)$
has $(q+1)(q^4+q^2+1)$ points and as many lines. By Feit and Higman \cite{FeitHigman}, thick generalized $n$-gons exist
only for $n\in\{2,3,4,6,8\}$, and a generalized octagon of order $(s,s)$ does not exist; so $\Mb(\Delta,D)$ is attained,
for $\Delta\ge3$, only at $D\in\{2,3,4,6\}$.

\subsubsection*{The projective plane}
Points of $\PG(2,q)$ are written $(x:y:z)$ and lines $[u:v:w]$, with $(x:y:z)\I[u:v:w]$ if and only if $ux+vy+wz=0$.

\subsubsection*{Quadrangles}
$W(3,q)$ is the symplectic quadrangle: the points of $\PG(3,q)$ and the lines of $\PG(3,q)$ that are totally isotropic
for a non-degenerate alternating form; it is a GQ of order $(q,q)$ \cite{PT}. $Q(4,q)$ is the parabolic quadric of
$\PG(4,q)$ with the lines it contains; it is a GQ of order $(q,q)$, dual to $W(3,q)$ \cite[3.2.1]{PT}. The
Ahrens--Szekeres quadrangle $\AS(q)$ ($q$ odd) is introduced in coordinates in Section~\ref{ssec:A}.

\subsubsection*{The split Cayley hexagon}\label{ssec:Hq}
Let $V=\F^7$ with coordinates $(X_0,\dots,X_6)$,
\[
Q(X)=X_0X_4+X_1X_5+X_2X_6-X_3^2,\qquad b(X,Y)=Q(X+Y)-Q(X)-Q(Y),
\]
and let $Q(6,q)$ be the parabolic quadric $\{\langle x\rangle:Q(x)=0\}$ of $\PG(6,q)$. Points $\langle x\rangle,\langle
y\rangle$ of $Q(6,q)$ are \emph{perpendicular}, $x\perp y$, if $b(x,y)=0$, i.e.\ if the line $xy$ lies on $Q(6,q)$. The
\emph{split Cayley hexagon} $H(q)$ \cite{T59}, \cite[\S2.4]{VM98} has as points all
$N=N(q)=(q^6-1)/(q-1)=(q+1)(q^4+q^2+1)$ points of $Q(6,q)$, and as lines the lines of $Q(6,q)$ whose Grassmann
coordinates $p_{ij}=x_iy_j-x_jy_i$ satisfy Tits' equations
\begin{equation}\label{eq:tits}
p_{12}=p_{34},\quad p_{54}=p_{32},\quad p_{20}=p_{35},\quad p_{65}=p_{30},\quad p_{01}=p_{36},\quad p_{46}=p_{31}.
\end{equation}
For example, the $q+1$ lines through $e_0=(1,0,\dots,0)$ are $\langle e_0,\lambda e_5+\mu e_6\rangle$,
$(\lambda:\mu)\in\PG(1,q)$; they lie in the totally singular plane $\langle e_0,e_5,e_6\rangle$. The same equations
define $H(q)$ for $q$ even, with the quadric written $X_0X_4+X_1X_5+X_2X_6+X_3^2$ (Section~\ref{ssec:evenq}).
\todo{The referee of the diameter-7 note recalls that this form of the quadric and of Tits' equations is
\cite[\S2.4.13]{VM98}. We could not consult [VM98] to confirm the place; please check and cite it.}

\begin{fact}[\cite{T59}, {\cite[\S2.4]{VM98}}]\label{fact:hex}
$H(q)$ is a generalized hexagon of order $(q,q)$. The lines of $H(q)$ through a point $x$ are the lines through $x$ of a
totally singular plane $\pi_x\subset x^\perp$. The collineation group contains $G_2(q)$; for $A\in\SL(3,q)$ the map
$(v,a,w)\mapsto(Av,a,A^{-T}w)$, with $x=(v,a,w)$, $v=(x_0,x_1,x_2)$, $a=x_3$, $w=(x_4,x_5,x_6)$, is a collineation. $H(q)$
has a duality (an incidence-preserving bijection exchanging points and lines) if and only if $q$ is a power of $3$.
\end{fact}

Let $I=I(q)$ be the incidence graph of $H(q)$, with $2N$ vertices, and $\Gamma_i(x)=\{y:d(x,y)=i\}$. Two elements are
\emph{opposite} if they are at distance 6, and $\opp(x)=\Gamma_6(x)$. The following facts are standard; we include the
short proofs because Sections~\ref{sec:d7} and~\ref{sec:d8} use them repeatedly.

\begin{lemma}\label{lem:hexfacts}
Let $x,y$ be vertices of $I$.
\begin{enumerate}[label=(\alph*)]
\item If $d(x,y)\le5$ there is a unique geodesic from $x$ to $y$. If $1\le d(x,y)=k\le5$, exactly one neighbour of $y$ is
  at distance $k-1$ from $x$ and the other $q$ at distance $k+1$; if $k=6$ all $q+1$ neighbours of $y$ are at distance
  5. Hence $|\Gamma_i(x)|=(q+1)q^{i-1}$ for $1\le i\le5$ and $|\Gamma_6(x)|=q^5$.
\item \emph{(covering)} There is $z\in\Gamma_2(x)$ with $d(z,y)\le3$ if $x,y$ have different types, and with
  $d(z,y)\le4$ if they have the same type.
\item \emph{(orthogonality)} For points $x,y$: $d(x,y)\le4$ if and only if $x\perp y$. So $\opp(x)\cap\cP=\cP\setminus
  x^\perp$.
\end{enumerate}
\end{lemma}

\begin{proof}
(a) Two geodesics would close a cycle of length at most $10<12$. Neighbours of $y$ are at distance $k\pm1$ from $x$;
the predecessor on the geodesic is at $k-1$, and a second one would give a second geodesic; at $k=6$ the diameter
forces $5$. The counts follow by induction. (b) If $2\le d(x,y)\le6$ take the vertex at distance 2 from $x$ on a
geodesic to $y$; if $d(x,y)=1$ take $z\in\Gamma_1(y)\setminus\{x\}$; if $x=y$ any $z\in\Gamma_2(x)$. The largest value
of $d(z,y)$ is $6-2=4$ (same type) or $5-2=3$ (different types). (c) If $d(x,y)\le4$ there is a point $z$ with $x,y$
on lines of $H(q)$ through $z$ ($z=x$ if $d(x,y)\le2$); these lines lie in $\pi_z$, so $b(x,y)=0$. Conversely
$x^\perp\cap Q(6,q)$ is a cone with vertex $x$ over a $Q(4,q)$ and has $1+q(q^3+q^2+q+1)$ points, which by (a) is the
number $1+(q+1)q+(q+1)q^3$ of points at distance at most 4 from $x$.
\end{proof}

\subsection{Doubling with levels, and the hat lemma}\label{ssec:doubling}
A \emph{duplicate} of a vertex (a new vertex with the same neighbours) does not increase the diameter (see
\cite{DF84}, which we have not seen, and \cite[\S7]{D85a}). Delorme \cite[\S7]{D85a} refined this into a doubling with
levels, stated there without proof, and used it for the 2013 values at $D=4$ and $D=6$ \cite[\S8(a),(b)]{D85a}. We
state it for incidence structures and give a proof.

Let $\cA=(\cA_0\supseteq\cA_1\supseteq\dots\supseteq\cA_r\supseteq\cA_{r+1}=\emptyset)$ be a filtration of a set
$\cA_0$ of elements (points and lines may be mixed). The \emph{level} of $x$ is $\lambda(x)=\max\{k:x\in\cA_k\}$, and
$\lambda(x)=-1$ if $x\notin\cA_0$. Fix $m\ge2$. Two words over $[m]=\{1,\dots,m\}$ are \emph{comparable} if one is a
prefix of the other.

\begin{definition}[Delorme]\label{def:doubling}
The $m$-fold doubling $\cS^{(m)}_\cA$ has the \emph{versions} $x^w$, for $x$ an element of $\cS$ and
$w\in[m]^{\lambda(x)+1}$ ($x^w$ is a point if $x$ is a point and a line if $x$ is a line), and
\[
x^w\I y^u\iff x\I y\ \text{and}\ w,u\ \text{are comparable}.
\]
\end{definition}

So an element of level $-1$ has one version (identified with $x$), adjacent to every version of each of its neighbours;
and a version $x^w$ of level $k\ge0$ is adjacent, for each neighbour $y$ of $x$, to the single version
$y^{w_1\cdots w_{j+1}}$ if $\lambda(y)=j\le k$, and to the $m^{j-k}$ versions $y^{wu}$ ($|u|=j-k$) if $\lambda(y)=j>k$.
This is the exact neighbourhood of every version. The map $x^w\mapsto x$ preserves incidence, so
$d(x^w,y^u)\ge d(x,y)$.

The conditions on $\cA$ are (with $D\ge2$):
\begin{itemize}
\item[(L0)] every element outside $\cA_0$ is incident with at most one element of $\cA_0$;
\item[(L1)] for every $i\ge0$: every element incident with an element of $\cA_{i+1}$ lies in $\cA_i$, and every element
  of $\cA_i\setminus\cA_{i+1}$ is incident with at most one element of $\cA_{i+1}$;
\item[(C$_D$)] for any two elements $x,y$ (possibly $x=y$) in the same connected component of the substructure induced
  on $\cA_0$, there is $z\notin\cA_0$ with $d(x,z)+d(z,y)\le D$.
\end{itemize}

\begin{theorem}[Delorme {\cite[\S7]{D85a}}, (a)--(c)]\label{thm:delorme}
Let $\cS$ have diameter at most $D\ge2$ and let $m\ge2$.
\begin{enumerate}[label=(\alph*)]
\item If (C$_D$) holds, $\cS^{(m)}_\cA$ has diameter at most $D$.
\item If (L0) and (L1) hold, every version satisfies $\deg x^w\le\deg x+m-1$, with equality only if $x$ is incident
  with an element of level $\lambda(x)+1$.
\item $\cS^{(m)}_\cA$ has $|\cP|+(m-1)\sum_{k\ge0}m^k|\cA_k\cap\cP|$ points and
  $|\cL|+(m-1)\sum_{k\ge0}m^k|\cA_k\cap\cL|$ lines.
\item Conversely, if $\cS^{(m)}_\cA$ has diameter at most $D$, the conclusion of (C$_D$) holds for every pair
  $x,y\in\cA_0$.
\end{enumerate}
\end{theorem}

Delorme states (a)--(c) without proof (``je ne donne pas la preuve''); (d) is not in \cite{D85a}.

\begin{proof}
(c) An element of level $k$ has $m^{k+1}=1+(m-1)(1+m+\dots+m^k)$ versions.

(b) Let $\lambda(x)=k$. By the description after Definition~\ref{def:doubling}, only neighbours of level $>k$
contribute more than one version. If $k=-1$, (L0) gives at most one neighbour in $\cA_0$; it has level exactly 0 (a
neighbour in $\cA_1$ would put $x$ in $\cA_0$ by (L1)), and contributes $m$ versions. If $k\ge0$, (L1) with $i=k$ gives
at most one neighbour in $\cA_{k+1}$, and (L1) with $i=k+1$ gives none in $\cA_{k+2}$; that neighbour contributes $m$
versions. So $\deg x^w\le\deg x-1+m$.

(a) For an infinite word $\omega$ the \emph{sheet} $\Sigma_\omega$ is the set of versions $x^w$ with $w$ a prefix of
$\omega$. It contains one version of each element, and $x^w\mapsto x$ is an isomorphism from the substructure on
$\Sigma_\omega$ onto $\cS$. Two versions with comparable words lie in a common sheet, hence are at distance at most
$D$. Let $x^w,y^u$ have incomparable words; both are non-empty, so $x,y\in\cA_0$. If $x,y$ are in different components
of the substructure on $\cA_0$, a shortest path $x=z_0,\dots,z_k=y$ ($k\le D$) passes through some $z_i\notin\cA_0$,
whose unique version lies in every sheet. Lift $z_0,\dots,z_i$ into a sheet through $x^w$ and $z_i,\dots,z_k$ into a
sheet through $y^u$: this is a path of length $k$. If $x,y$ are in the same component (including $x=y$), (C$_D$) gives
$z\notin\cA_0$ with $d(x,z)+d(z,y)\le D$, and the same lifting gives a path of length at most $D$.

(d) Take versions $x^w,y^u$ with $w$ starting with the letter 1 and $u$ with the letter 2, and a path of length at most
$D$ between them. Consecutive versions on it have comparable words, and comparable non-empty words start with the same
letter; so some version on the path has the empty word, that is, it is an element $z\notin\cA_0$. Projecting the path
gives a walk $x\to z\to y$ of length at most $D$.
\end{proof}

\begin{lemma}[hat lemma]\label{lem:hat}
Let $\cS$ be a generalized $n$-gon of order $(s,t)$, $s,t\ge1$, and $\cA_0$ a set of elements. Then (C$_n$) holds if
and only if every connected component of the substructure induced on $\cA_0$ is a tree of diameter at most $n-2$.
\end{lemma}

\begin{proof}
Suppose a component $B$ contains a path $x=z_0,\dots,z_{n-1}=y$ of length $n-1$ (it does if $B$ has diameter at least
$n-1$, and if $B$ contains a cycle, since cycles have length at least $2n$). By the girth this is the only path of
length at most $n-1$ from $x$ to $y$, so $d(x,y)=n-1$. If $d(x,z)+d(z,y)\le n$ then by parity
$d(x,z)+d(z,y)=n-1$, so $z$ lies on a shortest $x$--$y$ path, i.e.\ $z=z_i\in\cA_0$. So (C$_n$) fails.

Conversely let $B$ be a tree of diameter at most $n-2$, $x,y\in B$, and $P$ the path from $x$ to $y$ in $B$, of length
$d$. The incidence graph is connected and has cycles, so some element of $B$ has a neighbour outside $\cA_0$. Let $k$
be the least distance in $B$ from $P$ to such an element $u$, and $z$ a neighbour of $u$ outside $\cA_0$. Then
$d(x,z)+d(z,y)\le d+2k+2$. If $k=0$ this is at most $n$. If $k\ge1$, every element of $B$ at distance less than $k$ from
$P$ has all its neighbours in $B$, and every element has at least two neighbours. So we can walk $k$ steps away from
$P$ starting at $x$, and $k$ steps away from $P$ starting at $y$ (along two different neighbours if $x=y$); in the tree
$B$ this gives a path of length $d+2k$. Hence $d+2k\le n-2$ and $d(x,z)+d(z,y)\le n$.
\end{proof}

\begin{corollary}\label{cor:hat}
In a generalized $n$-gon of order $(s,t)$, $s,t\ge1$, for any filtration and $m\ge2$, $\cS^{(m)}_\cA$ has diameter at
most $n$ if and only if every component of the substructure on $\cA_0$ is a tree of diameter at most $n-2$. For $n=3$
the allowed components are single elements and flags; for $n=4$ also \emph{stars} (a point with some lines on it, or a
line with some of its points); for $n=6$ trees of diameter at most 4.
\end{corollary}

\begin{proposition}[diameter 3 needs patches]\label{prop:patches}
In a projective plane, if $\cA$ satisfies (L0), (L1) and (C$_3$), then $\cA_0$ has at most one point and at most one
line, and $\cA_1=\emptyset$.
\end{proposition}

\begin{proof}
If $x\ne y$ are points of $\cA_0$, the line $xy$ is either outside $\cA_0$ and incident with two elements of $\cA_0$
(against (L0)), or in $\cA_0$, and then $x,xy,y$ is a path of length $2>n-2$ in one component (against
Lemma~\ref{lem:hat}). Lines likewise. An element of $\cA_1$ would have its $q+1\ge3$ neighbours in $\cA_0$ by (L1).
\end{proof}

So pure doubling gains at most $2(m-1)$ elements in a projective plane; the hats of Section~\ref{sec:d3} double flags
and then change a bounded number of incidences.

%======================================================================
\section{Diameter 3: the fourth hat}\label{sec:d3}
%======================================================================

Araujo-Pardo and L\'opez \cite{AL20} add three \emph{hats} to $\PG(2,q)$ and obtain $\nb(q+2,3)\ge2(q^2+q+5)$. We add
a fourth. By Proposition~\ref{prop:patches}, every hat must change some incidences of the plane; the fourth hat costs
one rotation of three incidences.

\subsection{Frame and notation}
Let $\ell=[1:0:0]$ (the line $x=0$) and $A=(1:0:0)\notin\ell$. The points of $\ell$ and the lines on $A$ are
\[
X_t=(0:t:1),\quad P=(0:1:0);\qquad \lambda_t=[0:1:-t],\quad \lambda_\infty=[0:0:1]\qquad(t\in\F).
\]
Every point is listed once among $A$, $X_t$ ($t\in\F$), $P$, $(x:y:1)$ ($x\ne0$), $(x:1:0)$ ($x\ne0$); every line
once among $\ell$, $\lambda_t$, $\lambda_\infty$, $[u:1:w]$ ($u\ne0$), $[1:0:w]$ ($w\ne0$). We use these facts
(one-line computations):
\begin{itemize}
\item[(F1)] $\lambda_t\cap\ell=X_t$, $\lambda_\infty\cap\ell=P$; the points of $\lambda_t$ are $A$ and $(x:t:1)$,
  $x\in\F$ ($X_t$ at $x=0$); those of $\lambda_\infty$ are $A$ and $(x:1:0)$, $x\in\F$.
\item[(F2)] The lines on $X_t$ are $\ell$ and $[w:1:-t]$, $w\in\F$ ($\lambda_t$ at $w=0$); the lines on $P$ are $\ell$,
  $\lambda_\infty$ and $[1:0:w]$, $w\ne0$.
\item[(F3)] If $u=(x:y:z)$ with $x\ne0$, then $u$ and $X_s$ lie on $[(sz-y)/x:1:-s]\ne\ell$.
\item[(F4)] A line $[u:v:w]$ with $u\ne0$ meets $\lambda_s$ in $(-(vs+w)/u:s:1)\ne A$; $[u:1:w]$ meets $\ell$ in
  $X_{-w}$.
\end{itemize}
Let $q\ge4$, let $a,c\in\F\setminus\{0,1\}$ be distinct, and put
\[
S=\{0,1,a,c\},\qquad T=\{0,1,c\},\qquad \sigma=(0\ c\ 1):\ 0\mapsto c\mapsto1\mapsto0 .
\]

\subsection{The construction}

\begin{definition}\label{def:hat4}
$\cS(q;a,c)$ has the points of $\PG(2,q)$ and five new points $Q,H_0,H_1,H_a,H_c$, the lines of $\PG(2,q)$ and five
new lines $M,K_0,K_1,K_a,K_c$, and exactly these incidences:
\begin{itemize}
\item[(C1)] every incidence of $\PG(2,q)$ except $X_s\I\lambda_s$ ($s\in S$), and the \emph{rotation}
  $X_0\I\lambda_c$, $X_c\I\lambda_1$, $X_1\I\lambda_0$ (that is, $X_t\I\lambda_{\sigma(t)}$ for $t\in T$);
\item[(C2)] $H_s$ lies on $[w:1:-s]$ for every $w\in\F$ ($s\in S$): the lines on $X_s$ other than $\ell$;
\item[(C3)] $K_s$ contains $(x:s:1)$ for every $x\in\F$ ($s\in S$): the points of $\lambda_s$ other than $A$;
\item[(C4)] $M$ contains $P$, $X_t$ ($t\in\F\setminus T$), and $H_0,H_1,H_a,H_c$;
\item[(C5)] $Q$ lies on $\lambda_\infty$, $\lambda_t$ ($t\in\F\setminus T$), and $K_0,K_1,K_a,K_c$;
\item[(C6)] the \emph{links} $A\I K_a$, $H_a\I\ell$, $H_0\I K_c$, $H_c\I K_1$, $H_1\I K_0$ (that is, $H_t\I
  K_{\sigma(t)}$ for $t\in T$).
\end{itemize}
$G(q;a,c)$ is its incidence graph. $R(q;a,c)$ is $G(q;a,c)$ together with the \emph{filler} edges
\[
(x:y:1)\sim[x:1:-y]\quad(x\ne0,\ y\notin S),\qquad (x:1:0)\sim[1:0:-1/x]\quad(x\ne0).
\]
\end{definition}

The hat pair $(H_s,K_s)$ is a copy of the flag $(X_s,\lambda_s)$ with the flag ``crossed'', and $(Q,M)$ is a copy of the
anti-flag $(A,\ell)$.

\begin{theorem}\label{thm:hat4}
For every prime power $q\ge4$ and distinct $a,c\in\F\setminus\{0,1\}$, the graphs $G(q;a,c)$ and $R(q;a,c)$ are
bipartite, have $2(q^2+q+6)$ vertices, maximum degree $q+2$ and diameter 3, and $R(q;a,c)$ is $(q+2)$-regular. Hence
\[
\nb(q+2,3)\ge2(q^2+q+6),
\]
attained by a regular graph.
\end{theorem}

The construction and the proof use only the axioms of a projective plane together with an anti-flag $(A,\ell)$ and four
points $X_s$ of $\ell$, so a coordinate-free version holds in every projective plane of order $q\ge3$ (in coordinates
the four points are labelled by $0,1,a,c$, which needs $q\ge4$). It also has a regular completion there. At $q=3$ the
four points $X_s$ are all of $\ell$, there is no deficient element, and the structure is already 5-regular. For $q\ge4$
put $N_0=(q-1)(q-3)$, the number of points of degree $q+1$ and of lines of degree $q+1$ (see the table below). A
deficient point $p$ lies on exactly $q-4$ deficient lines: the lines on $p$ other than $Ap$ meet $\ell$ in the $q$
points other than $Ap\cap\ell$, four of which are the $X_s$. Dually, a deficient line has exactly $q-4$ deficient
points. So the bipartite graph of non-incident pairs between deficient points and deficient lines is regular of degree
$N_0-(q-4)=q^2-5q+7\ge1$, and a regular bipartite graph of positive degree has a perfect matching (K\H{o}nig); adding
it as flags completes the structure. In coordinates, the filler of Definition~\ref{def:hat4} is such a matching.

\subsection{Every neighbourhood}
Table~\ref{tab:hat4} lists the neighbourhood in $G(q;a,c)$ of every vertex, new or old. Here $\sigma^{-1}(0)=1$,
$\sigma^{-1}(1)=c$, $\sigma^{-1}(c)=0$.

\begin{longtable}{@{}l p{8.4cm} l@{}}
\caption{The neighbourhood of every vertex of $G(q;a,c)$.}\label{tab:hat4}\\
\toprule vertex & neighbours & degree\\ \midrule\endfirsthead
\toprule vertex & neighbours & degree\\ \midrule\endhead
\bottomrule\endfoot
$Q$ & $\lambda_\infty$, $\lambda_t$ ($t\notin T$), $K_0,K_1,K_a,K_c$ & $q+2$\\
$H_a$ & $[w:1:-a]$ ($w\in\F$), $M$, $\ell$ & $q+2$\\
$H_t$, $t\in T$ & $[w:1:-t]$ ($w\in\F$), $M$, $K_{\sigma(t)}$ & $q+2$\\
$M$ & $P$, $X_t$ ($t\notin T$), $H_0,H_1,H_a,H_c$ & $q+2$\\
$K_a$ & $(x:a:1)$ ($x\in\F$), $Q$, $A$ & $q+2$\\
$K_t$, $t\in T$ & $(x:t:1)$ ($x\in\F$), $Q$, $H_{\sigma^{-1}(t)}$ & $q+2$\\
$A$ & $\lambda_t$ ($t\in\F$), $\lambda_\infty$, $K_a$ & $q+2$\\
$P$ & $\ell$, $\lambda_\infty$, $[1:0:w]$ ($w\ne0$), $M$ & $q+2$\\
$X_t$, $t\notin S$ & $\ell$, $[w:1:-t]$ ($w\in\F$), $M$ & $q+2$\\
$X_a$ & $\ell$, $[w:1:-a]$ ($w\ne0$), $M$, $K_a$ & $q+2$\\
$X_t$, $t\in T$ & $\ell$, $[w:1:-t]$ ($w\ne0$), $\lambda_{\sigma(t)}$, $K_t$ & $q+2$\\
$(x:y:1)$, $x\ne0$ & $[u:1:-(ux+y)]$ ($u\in\F$), $[1:0:-x]$, and $K_y$ if $y\in S$ & $q+2$ if $y\in S$, else $q+1$\\
$(x:1:0)$, $x\ne0$ & $\lambda_\infty$, $[1:-x:w]$ ($w\in\F$) & $q+1$\\
$\ell$ & $P$, $X_t$ ($t\in\F$), $H_a$ & $q+2$\\
$\lambda_\infty$ & $A$, $(x:1:0)$ ($x\in\F$), $Q$ & $q+2$\\
$\lambda_t$, $t\notin S$ & $A$, $(x:t:1)$ ($x\in\F$), $Q$ & $q+2$\\
$\lambda_a$ & $A$, $(x:a:1)$ ($x\ne0$), $Q$, $H_a$ & $q+2$\\
$\lambda_t$, $t\in T$ & $A$, $(x:t:1)$ ($x\ne0$), $X_{\sigma^{-1}(t)}$, $H_t$ & $q+2$\\
$[u:1:w]$, $u\ne0$ & $(x:-ux-w:1)$ ($x\in\F$), $(-1/u:1:0)$, and $H_{-w}$ if $-w\in S$ & $q+2$ if $-w\in S$, else $q+1$\\
$[1:0:w]$, $w\ne0$ & $(-w:y:1)$ ($y\in\F$), $P$ & $q+1$\\
\end{longtable}

The vertices of degree $q+1$ are the points $(x:y:1)$ ($x\ne0$, $y\notin S$) and $(x:1:0)$ ($x\ne0$), and the lines
$[u:1:w]$ ($u\ne0$, $-w\notin S$) and $[1:0:w]$ ($w\ne0$): $(q-1)(q-3)$ of each kind. In $R(q;a,c)$ each of them gains
exactly one neighbour, its filler partner.

\subsection{Proof of Theorem~\ref{thm:hat4}}

\emph{Degrees.} Each row of Table~\ref{tab:hat4} is (F1) or (F2) corrected by (C1)--(C6), and within a row the listed
neighbours are distinct. For instance: $A$ keeps its $q+1$ lines and gains only $K_a$ ($A$ is not of the form
$(x:s:1)$); $X_a$ loses $\lambda_a$ and gains $M$ (as $a\notin T$) and $K_a$; $X_t$ ($t\in T$) loses $\lambda_t$, gains
$\lambda_{\sigma(t)}$ and $K_t$, and is not on $M$; a point $(x:y:1)$ with $x\ne0$ is not on $\ell$, so it loses
nothing, and it gains $K_s$ if and only if $y=s\in S$; $Q$ is on $1+(q-3)+4=q+2$ lines; $H_s$ is on $q$ lines by (C2),
on $M$, and on exactly one link line. The line rows are the dual count: $\ell$ gains only $H_a$; $\lambda_a$ loses
$X_a$ and gains $Q,H_a$; $\lambda_t$ ($t\in T$) loses $X_t$ and gains $X_{\sigma^{-1}(t)}$ and $H_t$; $[u:1:w]$ ($u\ne0$)
gains $H_s$ if and only if $s=-w\in S$; $M$ has $1+(q-3)+4$ points; $K_s$ has $q$ points, $Q$, and one link point.

The filler map $\delta:(x:y:z)\mapsto[x:z:-y]$ is a correlation of $\PG(2,q)$ (since $ux+wz+(-v)(-y)=ux+vy+wz$), a point
lies on its image if and only if $x^2=0$, i.e.\ if and only if it is on $\ell$, and $\delta$ maps $(x:y:1)\mapsto[x:1:-y]$
and $(x:1:0)\mapsto[x:0:-1]=[1:0:-1/x]$. So the filler is a bijection from the deficient points to the deficient lines,
and its pairs are anti-flags. Hence $G(q;a,c)$ has maximum degree $q+2$ and $R(q;a,c)$ is $(q+2)$-regular.

\emph{Reduction.} By Lemma~\ref{lem:diamgeom}(ii) it suffices to show that any two points are collinear and any two lines
concurrent. In the next two lemmas ``incident'' means incident in $\cS(q;a,c)$. The only incidences of $\PG(2,q)$ that are
missing are $X_s\I\lambda_s$, $s\in S$; so an incidence $u\I L$ of the plane survives whenever $u\notin\ell$ or $L$ is
not on $A$.

\begin{lemma}\label{lem:hat4points}
Any two points $u\ne v$ of $\cS(q;a,c)$ are collinear.
\end{lemma}

\begin{proof}
(P1) Both in $\PG(2,q)$; let $L=uv$. If both incidences survive, use $L$. Otherwise, say $u=X_s$ and $L=\lambda_s$ with
$s\in S$, and $v=A$ or $v=(x:s:1)$ with $x\ne0$. If $v=(x:s:1)$: use $K_s$. If $v=A$ and $s=a$: $K_a$ (by (C3), (C6)). If
$v=A$ and $s\in T$: $\lambda_{\sigma(s)}$ (rotation).

(P2) $Q$ and $u$. If $u\in\{A,P\}$ or $u=(x:1:0)$: $\lambda_\infty$. If $u=X_t$, $t\notin S$: $\lambda_t$. If $u=X_s$, $s\in
S$: $K_s$. If $u=(x:y:1)$, $x\ne0$: $\lambda_y$ if $y\notin T$, and $K_y$ if $y\in T$.

(P3) $H_s$ and $u$. If $u\notin\ell$ (including $u=A$): the line of (F3), which contains $H_s$ by (C2). If $u=X_s$:
$[1:1:-s]$. If $u=P$, or $u=X_t$ with $t\notin T$, $t\ne s$: $M$. If $u=X_t$ with $t\in T\setminus\{s\}$ and $s=a$: $\ell$.
If $u=X_t$ with $t\in T\setminus\{s\}$ and $s\in T$: as $\sigma$ is a 3-cycle, $t=\sigma(s)$ or $t=\sigma^{-1}(s)$; use
$K_{\sigma(s)}$ (it contains $X_{\sigma(s)}$ and $H_s$) or $\lambda_s$ (it contains $X_{\sigma^{-1}(s)}$ by the rotation, and
$H_s$ by (C2) with $w=0$).

(P4) Two new points. $H_s,H_{s'}$: $M$. $Q,H_a$: $\lambda_a$. $Q,H_t$ ($t\in T$): $K_{\sigma(t)}$.
\end{proof}

\begin{lemma}\label{lem:hat4lines}
Any two lines $L\ne N$ of $\cS(q;a,c)$ are concurrent.
\end{lemma}

\begin{proof}
(N1) Both in $\PG(2,q)$; let $p=L\cap N$. If both incidences survive, use $p$. Otherwise, say $p=X_s$ and $L=\lambda_s$,
and $N=\ell$ or $N=[w:1:-s]$ with $w\ne0$. If $N=[w:1:-s]$: $H_s$. If $N=\ell$ and $s=a$: $H_a$. If $N=\ell$ and $s\in T$:
$X_{\sigma^{-1}(s)}$ (rotation).

(N2) $M$ and $L$. If $L$ is on $P$ ($L=\ell$, $\lambda_\infty$ or $[1:0:w]$): $P$. If $L=\lambda_t$, $t\notin S$: $X_t$. If
$L=\lambda_s$, $s\in S$: $H_s$. If $L=[u:1:w]$, $u\ne0$: it meets $\ell$ in $X_r$, $r=-w$ (F4); use $X_r$ if $r\notin T$,
and $H_r$ if $r\in T$ ($H_r$ is on $[u:1:-r]$ and on $M$).

(N3) $K_s$ and $L$. If $L$ is not on $A$ ($L=\ell$ or $L=[u:v:w]$ with $u\ne0$): $L\cap\lambda_s=(x:s:1)$ (F1, F4), which
is on $K_s$, and its incidence with $L$ survives. If $L=\lambda_s$: $(1:s:1)$. If $L=\lambda_\infty$ or $L=\lambda_t$,
$t\notin T$: $Q$. If $L=\lambda_t$ with $t\in T\setminus\{s\}$ and $s=a$: $A$. If $L=\lambda_t$ with $t\in T\setminus\{s\}$
and $s\in T$: if $t=\sigma(s)$, use $X_s$ (on $\lambda_{\sigma(s)}$ and $K_s$); if $t=\sigma^{-1}(s)$, use $H_t$ (on
$\lambda_t$ and on $K_{\sigma(t)}=K_s$).

(N4) Two new lines. $K_s,K_{s'}$: $Q$. $M,K_a$: $X_a$. $M,K_t$ ($t\in T$): $H_{\sigma^{-1}(t)}$.
\end{proof}

\begin{proof}[Proof of Theorem~\ref{thm:hat4}]
By Lemmas~\ref{lem:hat4points} and~\ref{lem:hat4lines} and Lemma~\ref{lem:diamgeom}(ii), $G(q;a,c)$ has diameter at
most 3, and adding the filler does not increase it (Lemma~\ref{lem:flags}). The graphs are bipartite and not complete
bipartite ($A\notin\ell$), so the diameter is exactly 3. The order is $2(q^2+q+1)+10$.
\end{proof}

The proof uses only that $\F$ is a field and that $0,1,a,c$ are distinct. The correlation $\delta$, extended by
$H_s\leftrightarrow K_s$ and $Q\leftrightarrow M$, maps $\cS(q;a,c)$ built with $\sigma$ onto the dual of the structure
built with $\sigma^{-1}$, so Lemma~\ref{lem:hat4lines} is also Lemma~\ref{lem:hat4points} for $\sigma^{-1}$; we wrote both
out so that every case can be checked directly. Both the rotation and the link $H_1\I K_0$ are needed: without either,
the diameter is 4 in every case we checked by computer (every prime power $q\le13$ and every admissible $(a,c)$).

\begin{remark}[the three hats of {\cite{AL20}} in these coordinates]\label{rem:AL}
The structure of \cite{AL20} has hats $\{0,1,a\}$, no rotation, only the flag $(X_a,\lambda_a)$ removed, links $A\I
K_a$, $H_a\I\ell$, $H_0\I K_1$, $H_1\I K_0$; $M$ contains $P$ and $X_t$ ($t\ne0,1$), and $Q$ lies on $\lambda_\infty$ and
$\lambda_t$ ($t\ne0,1$). The dictionary with the affine coordinates of \cite{AL20} is $(x,y)=(x:y:1)$, $[m,b]=[m:-1:b]$,
$L_c=[1:0:-c]$, $L=[0:0:1]$, $P_m=(1:m:0)$; so $L_0=\ell$, $P_0=A$, $(0,s)=X_s$, $[0,s]=\lambda_s$ and $L=\lambda_\infty$.
In the completing matching of \cite[arXiv v1, p.~8]{AL20} the index set of $P_i\sim L_i$ should be $\F^*$; with
$\F\setminus\{0,1,a\}$ the four vertices $P_1,P_a,L_1,L_a$ keep degree $q+1$. Neither the order nor the diameter is
affected.
\end{remark}
\todo{Gabriela, Nacho: please check Remark~\ref{rem:AL} against the journal version of [AL20] (we have only seen arXiv
v1), and decide whether the remark should stay in the paper.}

\subsection{Cells}
For $q=9,11,13$, Theorem~\ref{thm:hat4} gives $(11,3)$: 192 (2013 table: 190), $(13,3)$: 276 (2013: 270; current value
274 \cite{AL20}) and $(15,3)$: 376 (2013: 370; current value 374 \cite{AL20}), all with regular graphs. (The 2013 value
190 at $(11,3)$ is due to Feria-Pur\'on, Miller and Pineda-Villavicencio \cite{FMP13} and, independently, to
\cite{AL20}.) At $q=5$ it gives $72<80=\nb(7,3)$, a value known to be exact \cite{FMP13}.

%======================================================================
\section{Diameter 4: clones and shadow lines in generalized quadrangles}\label{sec:d4}
%======================================================================

\subsection{Clones and shadow lines}
Copies of points and lines of a GQ go back to Delorme \cite[\S\S7--8]{D85a}: his doubling may copy any induced
substructure, points and lines together, and his $D=4$ construction (Example~\ref{ex:delorme4}) doubles a spread.
Araujo-Pardo, Dalf\'o, Fiol and L\'opez \cite[Sec.~4]{ADFL21} double every line of a GQ to get biregular graphs. The
construction below is not a doubling: a clone is joined to every old line through its point, so no tree condition on the
copied set is needed, and it allows weights larger than one, repeated lines, and ``stand-ins''. (A level-0 line all of
whose points are at level $-1$ is a shadow line with trivial stand-ins.)

\begin{definition}[clones and shadow lines]\label{def:clones}
Let $\cS=(\cP,\cL)$ be a GQ of order $(s,t)$, let $\omega:\cP\to\Z_{\ge0}$ be a \emph{weight}, and $\cM$ a finite
multiset of lines. For each entry $\mu$ of $\cM$, with underlying line $L_\mu$, choose for every point $p\I L_\mu$ a
\emph{stand-in} $\theta_\mu(p)\in\{p,p^1,\dots,p^{\omega(p)}\}$. The incidence structure $\cS[\omega,\cM,\theta]$ has
\begin{itemize}
\item points: $\cP$ and the \emph{clones} $p^1,\dots,p^{\omega(p)}$ of each $p\in\cP$;
\item lines: $\cL$ and one \emph{shadow line} $L^\mu$ for each entry $\mu$ of $\cM$;
\item incidences: those of $\cS$; $p^i\I L$ whenever $p\I L\in\cL$; and $\theta_\mu(p)\I L^\mu$ for $p\I L_\mu$.
\end{itemize}
\end{definition}

\begin{table}[ht]
\caption{Neighbourhoods in $\cS[\omega,\cM,\theta]$.}\label{tab:clones}
\begin{tabular}{@{}l p{7.6cm} l@{}}
\toprule vertex & neighbours & degree\\ \midrule
clone $p^i$ & the $t+1$ old lines on $p$; the shadow lines $L^\mu$ with $\theta_\mu(p)=p^i$ &
$t+1+\#\{\mu:\theta_\mu(p)=p^i\}$\\
shadow line $L^\mu$ & $\theta_\mu(p)$ for the $s+1$ points $p$ of $L_\mu$ & $s+1$\\
old point $p$ & its $t+1$ lines; the shadow lines $L^\mu$ with $\theta_\mu(p)=p$ & $t+1+\#\{\mu:\theta_\mu(p)=p\}$\\
old line $L$ & its $s+1$ points and all their clones & $s+1+\sum_{p\I L}\omega(p)$\\
\bottomrule
\end{tabular}
\end{table}

When all stand-ins are old points ($\theta=\mathrm{id}$), a shadow line is adjacent exactly to the old points of $L_\mu$,
and clones are adjacent only to old lines.

\begin{theorem}[clones and shadow lines]\label{thm:B}
Let $\cS$ be a GQ of order $(s,t)$, $s,t\ge1$, and $\cS'=\cS[\omega,\cM,\theta]$.
\begin{enumerate}[label=(\alph*)]
\item $\cS'$ has diameter 4.
\item $\cS'$ has maximum degree at most $\Delta$ if and only if $\Delta\ge\max(s,t)+1$, $\sum_{p\I L}\omega(p)\le\Delta-s-1$
  for every line $L$ of $\cS$, and every point of $\cS'$ (old point or clone) is the stand-in for at most $\Delta-t-1$
  entries of $\cM$.
\item $\cS'$ has $(s+1)(st+1)+\sum_p\omega(p)$ points and $(t+1)(st+1)+|\cM|$ lines.
\end{enumerate}
\end{theorem}

\begin{proof}
(b) and (c) are read off Table~\ref{tab:clones}. (a) Write $\pi$ for the map sending a clone or a shadow line to its
underlying point or line. Every point $x$ of $\cS'$ is incident with every old line through $\pi(x)$. We use
Lemma~\ref{lem:diamgeom}(i); no element is isolated. Let $(x,Y)$ be an anti-flag of $\cS'$.
\begin{itemize}
\item $Y$ old. Then $\pi(x)\notin Y$. Let $r$ be the point of $Y$ collinear with $\pi(x)$, on the line $N$. Then
  $x\I N\I r\I Y$.
\item $Y=L^\mu$ and $\pi(x)\notin L_\mu$. Let $r\I L_\mu$ be collinear with $\pi(x)$, on the line $N$. Then $x\I N\I
  \theta_\mu(r)\I L^\mu$.
\item $Y=L^\mu$ and $\pi(x)\I L_\mu$. Pick $r\I L_\mu$, $r\ne\pi(x)$ ($s\ge1$). Then $x\I L_\mu\I\theta_\mu(r)\I L^\mu$.
\end{itemize}
So $d(x,Y)\le3$ for every anti-flag, and the diameter is at most 4. Two non-collinear old points exist (a point is
collinear with $s(t+1)<(s+1)(st+1)-1$ others) and stay non-collinear: they share no old line, a shadow line $L^\mu$
contains only points over points of $L_\mu$, and an old point lies over itself only. So the diameter is 4.
\end{proof}

\begin{corollary}[exact weights, regular completion]\label{cor:exact}
Let $\Delta\ge\max(s,t)+1$, let $\omega$ be a weighted $(\Delta-s-1)$-ovoid (that is,
$\sum_{p\I L}\omega(p)=\Delta-s-1$ for every line $L$) and $\cM$ a weighted $(\Delta-t-1)$-cover (every point is on
exactly $\Delta-t-1$ entries of $\cM$). Then $\cS[\omega,\cM,\mathrm{id}]$ has diameter 4, maximum degree $\Delta$ and
$2\Delta(st+1)$ elements. It becomes $\Delta$-regular, with the same order and diameter 4, by the following explicit
flags: number the clones $c_0,c_1,\dots$ and the shadow lines $y_0,\dots,y_{b-1}$ ($b=|\cM|$), put $k=\Delta-t-1$, and
join $c_i$ to $y_j$ for $j\equiv ik,ik+1,\dots,ik+k-1\pmod b$. (If $b=0$ or $\sum_p\omega(p)=0$, the structure is already
$\Delta$-regular.)
\end{corollary}

\begin{proof}
Counting flags gives $\sum_p\omega(p)=(\Delta-s-1)(st+1)$ and $|\cM|=(\Delta-t-1)(st+1)$; the total is
$(s+t+2+\Delta-s-1+\Delta-t-1)(st+1)=2\Delta(st+1)$. Old points and lines have degree $\Delta$, clones $t+1$, shadow lines
$s+1$. The added pairs are anti-flags (with $\theta=\mathrm{id}$ no clone is on a shadow line); the blocks
$\{ik,\dots,ik+k-1\}$, $i=0,1,\dots$, run through $0,1,\dots,k\sum_p\omega(p)-1=b(\Delta-s-1)-1$, so modulo $b$ every
residue is hit exactly $\Delta-s-1$ times, and within a block the residues are distinct since $k\le b$. Adding flags
does not increase distances (Lemma~\ref{lem:flags}), and two non-collinear old points still share no line.
\end{proof}

\begin{example}[Delorme's $D=4$ construction {\cite[\S8(a)]{D85a}}; the 2013 values]\label{ex:delorme4}
In $W(3,q)$ take a spread $\Sigma$ (a set of $q^2+1$ lines partitioning the points), a line $v\in\Sigma$, and
$\cA_1=\{v\}$, $\cA_0=\Sigma\cup\{\text{points of }v\}$. (L1): the elements on $v$ are its points; a point of $v$ meets
only $v$ in $\cA_1$, and a line of $\Sigma\setminus\{v\}$ meets nothing in $\cA_1$. (L0): a point off $v$ is on exactly
one line of $\Sigma$ and on no other line of $\cA_0$; a line not in $\Sigma$ has at most one point on $v$. The
components of the substructure on $\cA_0$ are the star formed by $v$ and its points, and single lines, so (C$_4$) holds
(Lemma~\ref{lem:hat}). Theorem~\ref{thm:delorme} gives maximum degree $q+m$, diameter 4 and
$2(q+1)(q^2+1)+(m-1)\bigl((q^2+1)+(q+1)+m\bigr)$ elements. For $m=2$ this is $2q^3+3q^2+3q+6$: 346, 1734, 3064, 4946 at
$q=5,9,11,13$; for $m=3$, $q=13$ it is $4760+2\cdot187=5134$. These are the 2013 values at $\Delta=7,11,13,15,16$.
\end{example}

\begin{remark}[a limit for shadow lines on $Q(4,q)$]\label{rem:shadowlimit}
Let $q$ be odd, $\cO$ an ovoid of $Q(4,q)$ (a set of $q^2+1$ points meeting every line once), $\omega=\mathbf 1_\cO$,
and suppose the maximum degree is $q+2$. Then at most $q^2-q+2$ shadow lines are possible. Indeed every old point and
every clone serves at most one shadow line, so a point off $\cO$ is on at most one entry of $\cM$ and a point of $\cO$ on
at most two. Suppose some $o\in\cO$ is on two entries. They are distinct lines: if both were copies of one line $N$, each
of the $q$ points of $N$ other than $o$ is off $\cO$, has no clone, and would be the stand-in of both entries, giving it
degree $q+3$. Each other entry $L$ misses $o$ and so has a unique point $y$ collinear with $o$; $y\notin\cO$ (points of
an ovoid are pairwise non-collinear); the line $oy$ is not an entry of $\cM$ (otherwise $y\notin\cO$ would be on two
entries), so it is one of the $q-1$ lines on $o$ other than the two entries through $o$, and $y$ is one of its $q$ points
other than $o$; distinct $L$ give distinct $y$. So $|\cM|\le2+q(q-1)$. Otherwise every point is on at most one entry,
$\cM$ is a partial spread, and $|\cM|\le q^2-q+1$ (Tallini \cite{Tal91}; see also \cite[\S1]{CDPS22}).
\end{remark}

\begin{example}[$(7,4)\ge359$]\label{ex:359}
In $Q(4,5)$ take one clone of each point of an ovoid $\cO$ that is an elliptic quadric section ($\omega=\mathbf
1_\cO$), and 21 shadow lines with stand-ins. Theorem~\ref{thm:B} gives diameter 4, maximum degree 7, and
$2\cdot6\cdot26+26+21=359$ elements. Appendix~\ref{app:359} lists the 21 lines and the stand-ins in coordinates, so that
the conditions of Theorem~\ref{thm:B}(b) can be checked by hand. Every ovoid of $Q(4,p)$, $p$ prime, is an elliptic
quadric \cite{BGS06}, so the choice of ovoid does not matter. By Remark~\ref{rem:shadowlimit} at most $q^2-q+2=22$ shadow
lines are possible in this setting; 21 is the maximum for this family (an integer program, solved to optimality).
This example is not a member of a proved infinite family. The best regular graph we know at $(7,4)$ has 350 vertices
(Corollary~\ref{cor:A1} at $q=5$).
\end{example}
\todo{Example~\ref{ex:359}: the optimality of 21 comes from an integer-programming solver (HiGHS). Keep the sentence,
or drop it as computer-dependent?}

\subsection{The Ahrens--Szekeres quadrangle and two elliptic quadrics}\label{ssec:A}

\subsubsection*{Coordinates}
Let $q$ be odd. Let $W(q)=W(3,q)$ be the symplectic quadrangle of $\PG(3,q)$ for the alternating form
\[
\beta(u,v)=u_0v_1-u_1v_0+u_2v_3-u_3v_2,
\]
and let $x=(1:0:0:0)$; then $x^\perp$ is the plane $X_1=0$. Payne's derived quadrangle $\AS(q)=P(W(q),x)$
\cite{P71}, \cite[3.1.4]{PT} has
\begin{itemize}
\item \emph{points}: the $q^3$ points $(a:1:b:c)$, $a,b,c\in\F$ (the points off $x^\perp$); we write $(a,w)$ with
  $w=(b,c)$;
\item \emph{lines of type (i)} (\emph{fibres}), $q^2$ of them: $\{(a_0+\mu:1:w_0):\mu\in\F\}$, the line
  $\langle x,(a_0:1:w_0)\rangle$ minus $x$;
\item \emph{lines of type (ii)}, $q^2(q+1)$ of them: for a point $(a_0:1:w_0)$ and a direction $d=(d_1:d_2)\in\PG(1,q)$,
  \[
  \ell(a_0,w_0;d)=\{(a_0+\lambda\rho:1:w_0+\lambda d):\lambda\in\F\},\qquad \rho=\det(w_0,d)=b_0d_2-c_0d_1,
  \]
  the totally isotropic line through $(a_0:1:w_0)$ and $(\rho:0:d_1:d_2)$, minus the latter point.
\end{itemize}
(Check: $\beta((a_0,1,b_0,c_0),(\rho,0,d_1,d_2))=-\rho+b_0d_2-c_0d_1=0$. A point off $x^\perp$ is on $q+1$ totally isotropic
lines, none through $x$, each meeting $x^\perp$ in one point $(\rho:0:d)$, $d\ne0$; so each point of $\AS(q)$ is on one
fibre and on $q+1$ lines of type (ii), and every line has $q$ points.) $\AS(q)$ is a GQ of order $(q-1,q+1)$ \cite{P71},
\cite[3.1.4, 3.2.6]{PT}; for $q$ odd it is the Ahrens--Szekeres quadrangle \cite{AS69}. The fibres form a spread.

Fix a non-square $\eps\in\F$ and put $N(b,c)=b^2-\eps c^2$, an anisotropic form ($N(w)=0$ only for $w=0$). For
$\kappa\in\F^*$ let $\cE_\kappa$ be the quadric $X_0X_1=\kappa N(X_2,X_3)$. It is elliptic (a hyperbolic pair plus an
anisotropic binary form), passes through $x$ with tangent plane $x^\perp$, and
$\cE_\kappa\setminus\{x\}=\{(\kappa N(w):1:w):w\in\F^2\}$ is a set of $q^2$ points of $\AS(q)$, one on each fibre.

\begin{theorem}[a weighted 2-ovoid of $\AS(q)$]\label{thm:A}
Let $q$ be odd, $\eps$ a non-square and $\kappa_1\kappa_2=1/(4\eps)$. Every line of $\AS(q)$ meets
$\cE_{\kappa_1}\setminus\{x\}$ and $\cE_{\kappa_2}\setminus\{x\}$ in exactly two points in total (a point in both counted
twice). The two quadrics meet only in $x$ and $e=(0:1:0:0)$. So
$\varpi=\mathbf 1_{\cE_{\kappa_1}\setminus x}+\mathbf 1_{\cE_{\kappa_2}\setminus x}$ is a weighted 2-ovoid of $\AS(q)$, with
$\varpi(e)=2$ and $\varpi\in\{0,1\}$ elsewhere.
\end{theorem}

The choice $\kappa_1=\frac12$, $\kappa_2=\frac1{2\eps}$ gives the two quadrics
\[
\cE^{(1)}:\ 2X_0X_1=X_2^2-\eps X_3^2,\qquad \cE^{(2)}:\ 2\eps X_0X_1=X_2^2-\eps X_3^2 .
\]

\begin{corollary}\label{cor:A1}
For every odd prime power $q$, adjoining to $\AS(q)$ one clone of each point of $\cE^{(1)}\setminus\{x\}$ and one of each
point of $\cE^{(2)}\setminus\{x\}$ (so $e$ gets two clones) gives a $(q+2)$-regular incidence structure of diameter 4 with
$2q^2(q+2)$ elements. Hence $\nb(q+2,4)\ge2q^2(q+2)$, attained by a regular graph.
\end{corollary}

\emph{New vertices.} For $P=(\kappa_iN(w):1:w)\in\cE_{\kappa_i}\setminus\{x\}$ ($i=1,2$, $w\in\F^2$), the clone $P^{(i)}$ is
adjacent exactly to the $q+2$ lines of $\AS(q)$ through $P$: the fibre $\{(\kappa_iN(w)+\mu:1:w):\mu\in\F\}$ and, for each
$d\in\PG(1,q)$, the line $\ell(\kappa_iN(w),w;d)$. The two clones of $e$ are adjacent to the same $q+2$ lines. Changed old
neighbourhoods: the fibre over $w_0$ gains the clones of $(\kappa_1N(w_0):1:w_0)$ and $(\kappa_2N(w_0):1:w_0)$; the line
$\ell(a_0,w_0;d)$ gains the clones of its points on $\cE_{\kappa_1}$ and on $\cE_{\kappa_2}$, which correspond to the roots
$\lambda$ of the quadratic \eqref{eq:quadratic} below with $\kappa=\kappa_1,\kappa_2$ (two clones in total). Old points
are unchanged.

\begin{proof}[Proof of Theorem~\ref{thm:A}]
\emph{Fibres.} The fibre over $w_0$ meets $\cE_\kappa$ only at $\mu=\kappa N(w_0)-a_0$: one point of each quadric.

\emph{Lines of type (ii).} Let $B=N(w_0+d)-N(w_0)-N(d)=2(b_0d_1-\eps c_0d_2)$. The point with parameter $\lambda$ of
$\ell(a_0,w_0;d)$ lies on $\cE_\kappa$ if and only if $a_0+\lambda\rho=\kappa N(w_0+\lambda d)$, that is,
\begin{equation}\label{eq:quadratic}
\kappa N(d)\lambda^2+(\kappa B-\rho)\lambda+\kappa N(w_0)-a_0=0.
\end{equation}
As $N(d)\ne0$ this is a quadratic, with $1+\chi(\mathrm{disc}_\kappa)$ roots. A direct expansion gives
$B^2-4N(w_0)N(d)=4\eps(b_0d_2-c_0d_1)^2=4\eps\rho^2$, hence
\[
\mathrm{disc}_\kappa=(\kappa B-\rho)^2-4\kappa N(d)(\kappa N(w_0)-a_0)=4\kappa N(d)\,(a_0-a^*_\kappa),\qquad
a^*_\kappa=-\frac{4\eps\kappa\rho^2-2B\rho+\rho^2/\kappa}{4N(d)} .
\]
Then $a^*_{\kappa_1}-a^*_{\kappa_2}=-\frac{\rho^2}{4N(d)}(\kappa_1-\kappa_2)\bigl(4\eps-\frac1{\kappa_1\kappa_2}\bigr)=0$; call
the common value $a^*$. The line meets the two quadrics in
\[
\bigl(1+\chi(\kappa_1)\chi(N(d)(a_0-a^*))\bigr)+\bigl(1+\chi(\kappa_2)\chi(N(d)(a_0-a^*))\bigr)=2
\]
points, because $\chi(\kappa_1)\chi(\kappa_2)=\chi(1/(4\eps))=\chi(\eps)=-1$, so $\chi(\kappa_1)+\chi(\kappa_2)=0$. (If
$a_0=a^*$, the line is tangent to both quadrics and meets each in one point.)

\emph{Intersection.} $\kappa_1\ne\kappa_2$, since $\kappa_1^2=1/(4\eps)$ is impossible for a non-square $\eps$. So
$\kappa_1N(w)=\kappa_2N(w)$ forces $N(w)=0$, $w=0$, and the common points are $x$ and $e$.
\end{proof}

\begin{proof}[Proof of Corollary~\ref{cor:A1}]
Theorem~\ref{thm:B} and Corollary~\ref{cor:exact} with $(s,t)=(q-1,q+1)$, $\omega=\varpi$ of Theorem~\ref{thm:A} and
$\cM=\emptyset$: every line has $q$ points and gains exactly two clones, every old point is on $q+2$ lines, and every
clone is on the $q+2$ lines of its point. The order is $q^3+q^2(q+2)+2q^2=2q^2(q+2)$.
\end{proof}

\begin{remark}
Weights are necessary: $\AS(3)$ has no 1- or 2-ovoid and $\AS(5)$ has no $m$-ovoid for $m=1,\dots,4$ (computer-checked
by integer programming). No choice of clones and shadow lines on $\AS(q)$ does better at $\Delta=q+2$: every point of
$\AS(q)$ is already on $t+1=q+2$ lines, so no shadow line can be added, and each line has $q$ points and takes at most 2
clones, so $\sum_p\omega(p)\le2q^2$. Corollary~\ref{cor:A1} beats Delorme's $2q^3+3q^2+3q+6$ by $q^2-3q-6>0$ for $q\ge5$.
\end{remark}

\noindent\emph{Cells.} $(11,4)$: 1782 (2013: 1734); $(13,4)$: 3146 (3064); $(15,4)$: 5070 (4946). All regular.

\subsection{Singer 2-ovoids of $W(3,q)$}

\subsubsection*{The cyclic model}
Let $q$ be odd, $\F=\GF(q)\subset K=\GF(q^2)\subset E=\GF(q^4)$, and $g$ a primitive element of $E$. As $E$ is a
4-dimensional $\F$-space, $\PG(3,q)=E^*/\F^*$ (any $\F$-basis of $E$ gives homogeneous coordinates); the point $\langle
g^i\rangle$ has index $i$ modulo $(q+1)(q^2+1)$. Write $\bar y=y^{q^2}$ (the involution of $E/K$), $k'=k^q$ for $k\in K$,
$\Nm(y)=y\bar y\in K$ and $\Tr=\Tr_{K/\F}$. Let $d=g^{(q^2+1)/2}$; then $\bar d=-d$, and $h=d^2=g^{q^2+1}$ generates
$K^*$. The form
\[
b(y,z)=\Tr_{E/\F}(d\,y\bar z)=\Tr\bigl(d(y\bar z-\bar yz)\bigr)
\]
is alternating and non-degenerate; $W(3,q)$ is its symplectic quadrangle. Multiplication by $g^{q+1}$ multiplies $b$ by
$g^{(q+1)(q^2+1)}\in\F^*$, so it is a collineation of $W(3,q)$. Its point orbits are the classes
\[
C_a=\{\langle g^i\rangle:\ i\equiv a\pmod{q+1}\},\qquad |C_a|=q^2+1,
\]
and $\langle y\rangle\in C_a$ if and only if $\Nm(y)\in h^a\F^*$ (as $\Nm(g^i)=h^i$ and $K^*/\F^*$ is cyclic of order
$q+1$, generated by $h$). These are Ebert's elliptic quadrics partitioning $\PG(3,q)$ \cite{E85}.

The $K$-subspaces $yK$ ($y\in E^*$) are lines of $\PG(3,q)$ (the regular spread), with points
$\langle g^{i+(q^2+1)j}\rangle$, $j=0,\dots,q$. They are totally isotropic: for $k_1,k_2\in K$ we have
$\overline{yk_i}=\bar yk_i$ (as $\bar k=k^{q^2}=k$), so $b(yk_1,yk_2)=\Tr\bigl(d\,\Nm(y)(k_1k_2-k_1k_2)\bigr)=0$.
So $\Sigma=\{yK\}$ is a spread of $W(3,q)$ with $q^2+1$ lines.

\begin{proposition}[Singer 2-ovoids]\label{prop:C}
Let $q$ be odd.
\begin{enumerate}[label=(\alph*)]
\item Every line of $\PG(3,q)$ meets each class $C_a$ in at most two points.
\item Every totally isotropic line of $W(3,q)$ meets $C_a\cup C_{q-a}$ in exactly two points.
\end{enumerate}
So the $(q+1)/2$ sets $C_a\cup C_{q-a}$, $0\le a\le(q-1)/2$, are 2-ovoids of $W(3,q)$ partitioning its points. (The pairing
$a\leftrightarrow q-a$ has no fixed point, as $2a\equiv-1\pmod{q+1}$ is impossible for $q+1$ even.)
\end{proposition}

\begin{proof}
Let $L=\langle y,z\rangle_\F$ and $u=sy+tz$ ($s,t\in\F$). Then $\Nm(u)=A_1s^2+\tau st+A_2t^2$ with $A_1=\Nm(y)$,
$A_2=\Nm(z)$, $\tau=y\bar z+\bar yz$, all in $K$. Put $\theta=y\bar z-\bar yz$. As $\bar\theta=-\theta$, $\theta/d\in K$; write
$\theta=dk_0$ with $k_0\in K$. Then
\[
\mathfrak d:=\tau^2-4A_1A_2=\theta^2=hk_0^2,\qquad b(y,z)=\Tr(hk_0).
\]
Fix $\varkappa\in K^*$ with $\varkappa'=-\varkappa$. For $\mu\in K^*$ put $\nu=\varkappa/\mu$; the kernel of
$k\mapsto\Tr(\nu k)$ is the $\F$-subspace $\mu\F$. So $\langle u\rangle$ lies in
$\widetilde C_\mu:=\{\langle u\rangle:\Nm(u)\in\mu\F^*\}$ (thus $C_a=\widetilde C_{h^a}$) if and only if $(s:t)$ is a zero of
the binary form
\[
f_\nu(s,t)=\Tr(\nu A_1)s^2+\Tr(\nu\tau)st+\Tr(\nu A_2)t^2 .
\]
(a) Suppose $f_\nu\equiv0$. Then $A_1,\tau,A_2\in\mu\F$. If $\mathfrak d\ne0$, then $hk_0^2=\mathfrak d\in\mu^2\F^*$, so
$h\in(K^*)^2\F^*=(K^*)^2$ (every element of $\F^*$ is a square in $K$, since $q-1$ divides $(q^2-1)/2$); but $h$ generates
$K^*$, which has even order. If $\mathfrak d=0$, then $\Nm(sy+tz)=\mu(\alpha_0s^2+\beta_0st+\gamma_0t^2)$ with
$\alpha_0,\beta_0,\gamma_0\in\F$, $\beta_0^2=4\alpha_0\gamma_0$ and $\alpha_0\ne0$ (as $A_1\ne0$), which vanishes at
$(s:t)=(-\beta_0:2\alpha_0)$: impossible, since $\Nm(u)\ne0$ for $u\ne0$. So $f_\nu\not\equiv0$, and it has
$1+\chi(\mathrm{disc}(\nu))\le2$ projective zeros, where
$\mathrm{disc}(\nu)=\Tr(\nu\tau)^2-4\Tr(\nu A_1)\Tr(\nu A_2)$.

(b) Expanding with $\Tr(\nu k)=\nu k+\nu'k'$,
\[
\mathrm{disc}(\nu)=\Tr(\nu^2\mathfrak d)+\nu\nu'\Psi,\qquad \Psi=2\tau\tau'-4A_1A_2'-4A_1'A_2\in\F .
\]
If $L$ is totally isotropic, then $\Tr(hk_0)=0$, so $hk_0\in\varkappa\F$ and $\eta:=h\mathfrak d=(hk_0)^2\in\varkappa^2\F\subseteq\F$.
Then $\Tr(h^2\nu'^2\mathfrak d)=\eta(h\nu'^2+h'\nu^2)=hh'\,\Tr(\nu^2\mathfrak d)$, and since $(h\nu')(h\nu')'=hh'\nu\nu'$,
\[
\mathrm{disc}(h\nu')=\Nm_{K/\F}(h)\,\mathrm{disc}(\nu).
\]
$\Nm_{K/\F}(h)=h^{q+1}$ generates $\F^*$, so it is a non-square and $\chi(\mathrm{disc}(h\nu'))=-\chi(\mathrm{disc}(\nu))$. The
class belonging to $h\nu'$ is $\varkappa/(h\nu')=-\mu'/h\equiv(h\mu)^{-1}$ modulo $\F^*$ (using $\mu^{q+1}\in\F^*$); in
indices, $a\mapsto-a-1\equiv q-a\pmod{q+1}$. Scaling $\nu$ by $\F^*$ multiplies $\mathrm{disc}$ by a square. So $L$ meets
$C_a\cup C_{q-a}$ in $(1+\chi(\mathrm{disc}(\nu)))+(1-\chi(\mathrm{disc}(\nu)))=2$ points.
\end{proof}

\begin{remark}[literature]
For $q$ odd, $W(3,q)$ has $m$-ovoids for every even $m$ and $(q+1)/2$-ovoids \cite{BLP09}, and \cite{FWX19} states,
citing \cite{BLP09}, that $W(3,p^e)$ ($p$ odd) has a partition into 2-ovoids. Cossidente, Culbert, Ebert and Marino
\cite{CCEM08} construct $m$-ovoids of $W(3,q)$ with Singer-type stabilisers. We have not checked whether the 2-ovoids
$C_a\cup C_{q-a}$ coincide with those; Proposition~\ref{prop:C} is very likely known, and we claim only the short proof.
The proof uses only that $d^2$ is a non-square in $K$: for any $\delta_0$ with $\bar\delta_0=-\delta_0$, $C_a\cup C_b$ is a
2-ovoid of the form $\Tr_{E/\F}(\delta_0y\bar z)$ when $h^{a+b}\delta_0^2\in\F^*$.
\end{remark}
\todo{Proposition~\ref{prop:C}: do you know whether the 2-ovoids $C_a\cup C_{q-a}$ are those of [BLP09], [CCEM08] or
[FWX19]? If so, cite and shorten.}

\begin{corollary}\label{cor:C1}
Let $q$ be odd, $0\le a\le(q-1)/2$, $\cO=C_a\cup C_{q-a}$ and $\Sigma$ the regular spread $\{yK\}$. Adjoin to $W(3,q)$
\begin{itemize}
\item one clone $P'$ of each point $P\in\cO$, incident with the $q+1$ totally isotropic lines on $P$ (the lines $\langle
  P,Z\rangle$ with $b(P,Z)=0$);
\item two shadow lines $L',L''$ for each $L\in\Sigma$, each incident with the $q+1$ points of $L$.
\end{itemize}
The result is bipartite, has diameter 4, maximum degree $q+3$ (old points and lines have degree $q+3$, clones and shadow
lines $q+1$) and $2(q^2+1)(q+3)$ elements. Joining clone $c_i$ to shadow lines $y_{2i},y_{2i+1}$ (indices modulo
$2(q^2+1)$, any numbering) makes it $(q+3)$-regular with diameter 4. Hence $\nb(q+3,4)\ge2(q^2+1)(q+3)$, attained by a
regular graph.
\end{corollary}

\begin{proof}
$W(3,q)$ is a GQ of order $(q,q)$. By Proposition~\ref{prop:C}(b), $\omega=\mathbf 1_\cO$ is a 2-ovoid ($=\Delta-s-1$ for
$\Delta=q+3$), and $\cM=2\Sigma$ is a 2-cover (every point is on exactly one spread line). Apply
Corollary~\ref{cor:exact} with $k=2$.
\end{proof}

This is the point--line dual of clones of an ovoid of $Q(4,q)$ with a dual 2-cover; written in $W(3,q)$ it needs no
Klein correspondence. More generally, $j$ copies of the spread and $j/2$ copies of the 2-ovoid give
$\nb(q+1+j,4)\ge2(q^2+1)(q+1+j)$ for even $j$.

\begin{corollary}[all degrees]\label{cor:B2}
For every prime power $q$ and every $\Delta\ge q+1$ there is a $\Delta$-regular incidence structure of diameter 4 with
$2\Delta(q^2+1)$ elements if $q$ is even or $\Delta-q-1$ is even, and with $2\Delta q^2$ elements otherwise. In particular
$\nb(\Delta,4)\ge2q^2\Delta$ for every prime power $q\le\Delta-1$.
\end{corollary}

\begin{proof}
Put $j=\Delta-q-1$. In the first two cases we apply Corollary~\ref{cor:exact} in $W(3,q)$ ($s=t=q$, so
$\Delta-s-1=\Delta-t-1=j$). If $q$ is even, $W(3,q)$ has an ovoid $\cO_0$ and a spread $\Sigma$ \cite[3.4.1]{PT}; take
$\omega=j\mathbf 1_{\cO_0}$ and $\cM=j\Sigma$. If $q$ is odd and $j$ even: $\omega=(j/2)\mathbf 1_\cO$ with $\cO$ a 2-ovoid of
Proposition~\ref{prop:C}, and $\cM=j\Sigma$. If $q$ is odd and $j$ odd, then $\Delta\ge q+2$, and we apply
Corollary~\ref{cor:exact} in $\AS(q)$ ($s=q-1$, $t=q+1$, so $\Delta-s-1=\Delta-q$ and $\Delta-t-1=\Delta-q-2$), with
$(\Delta-q)/2$ copies of the weighted 2-ovoid of Theorem~\ref{thm:A} and $\Delta-q-2$ copies of the fibre spread.
\end{proof}

\noindent\emph{Cells.} $(16,4)$: $q=13$, $5440=2\cdot170\cdot16$ (2013: 5134), regular. In the range $\Delta\le16$ this is
the only open cell with $\Delta-3$ an odd prime power. Without stand-ins, 5440 is the most a structure of clones and
shadow lines on $W(3,13)$ or $Q(4,13)$ of maximum degree 16 can have (each line takes at most 2 clones and each point
lies on at most 2 shadow lines).

%======================================================================
\section{Diameter 6: whiskers in the split Cayley hexagon}\label{sec:d6}
%======================================================================

Throughout Sections~\ref{ssec:d6coords}--\ref{ssec:thmD}, $q$ is odd; even $q$ is treated in
Section~\ref{ssec:evenq}. We use the model of $H(q)$ in Section~\ref{sec:prelim}.

\subsection{Elliptic hyperplane sections and Hermitian spreads}\label{ssec:d6coords}
Let $v=\langle e_1,e_6\rangle$. It is a line of $Q(6,q)$ ($e_1,e_6$ are singular and $b(e_1,e_6)=0$) and a line of $H(q)$
(no equation of \eqref{eq:tits} involves $p_{16}$). Fix $c\in\F$ with $-c$ a non-square. For $z=(\alpha,\beta)\in\F^2$ let
\[
\Pi_z:\ X_0+cX_4+\alpha X_2+\beta X_5=0,\qquad \Pi=\Pi_{(0,0)},
\]
and let $S_z$ be the set of lines of $H(q)$ contained in $\Pi_z$, $S=S_{(0,0)}$. Every $\Pi_z$ contains $v$.

\begin{lemma}\label{lem:Piz}
For every $z\in\F^2$:
\begin{enumerate}[label=(\alph*)]
\item $\Pi_z\cap Q(6,q)$ is an elliptic quadric $Q^-(5,q)$.
\item Every point of $H(q)$ in $\Pi_z$ is on exactly one line of $S_z$. So $|S_z|=q^3+1$, the lines of $S_z$ are pairwise
  opposite, and $v\in S_z$ ($S_z$ is a Hermitian spread \cite{T80}).
\item A line of $H(q)$ not in $S_z$ meets $\Pi_z$ in exactly one point.
\item The points of $H(q)$ at distance at most 3 from $v$ are exactly the points of $Q(6,q)$ with $X_2=X_5=0$; among them
  only the points of $v$ lie in $\Pi_z$.
\item $S_z\cap S_{z'}=\{v\}$ for $z\ne z'$.
\end{enumerate}
\end{lemma}

\begin{proof}
(a) On $\Pi_z$ substitute $X_0=-cX_4-\alpha X_2-\beta X_5$:
$X_0X_4+X_1X_5+X_2X_6-X_3^2=(X_1-\beta X_4)X_5+X_2(X_6-\alpha X_4)-(X_3^2+cX_4^2)$, two hyperbolic planes plus the
anisotropic form $X_3^2+cX_4^2$. So the form on $\Pi_z$ is non-degenerate of Witt index 2.

(b) Let $y\in\Pi_z\cap Q(6,q)$. By (a), $\Pi_z\cap Q(6,q)$ contains no totally singular plane, so $\pi_y\not\subseteq\Pi_z$
and $\pi_y\cap\Pi_z$ is one line through $y$: the only line of $H(q)$ on $y$ inside $\Pi_z$. $Q^-(5,q)$ has
$(q+1)(q^3+1)$ points, so $|S_z|=q^3+1$. Two lines of $S_z$ are not concurrent (uniqueness). If two were at distance 4, a
line $N$ of $H(q)$ would meet both in distinct points, so $N\subseteq\Pi_z$, $N\in S_z$, concurrent with them: a
contradiction. So they are opposite.

(c) A line not in the hyperplane $\Pi_z$ meets it in one point.

(d) Let $y=ae_1+be_6$ be on $v$ and $yu$ a line of $H(q)$. Then $p_{12}=p_{34}$ gives $au_2=0$, $p_{65}=p_{30}$ gives
$bu_5=0$, and $b(y,u)=0$ gives $au_5+bu_2=0$; so $u_2=u_5=0$. Conversely $Q(6,q)\cap\{X_2=X_5=0\}$ is $X_0X_4=X_3^2$, a
cone with vertex $v$ over a conic, with $(q+1)+(q+1)q^2$ points: exactly the number of points at distance at most 3 from
$v$ ($q+1$ on $v$, and $q$ more on each of the $q(q+1)$ other lines meeting $v$). If such a point is in $\Pi_z$, then
$X_0=-cX_4$ and $X_3^2+cX_4^2=0$, so $X_0=X_3=X_4=0$ and the point is on $v$.

(e) If $\ell\in S_z\cap S_{z'}$ with $\ell\ne v$, then $(\alpha-\alpha')X_2+(\beta-\beta')X_5=0$ on $\ell$, so some point of
$\ell$ has $X_2=X_5=0$ and is at distance at most 3 from $v$ by (d); but $\ell$ and $v$ are opposite by (b), so every point
of $\ell$ is at distance 5 from $v$.
\end{proof}

\subsection{Delorme's levels and whiskers}
Let $\cP(v)$ be the points of $v$, $\cL_v$ the $q^2+q$ lines other than $v$ that meet $v$, and $\cP_\Pi$ the
$(q+1)(q^3+1)$ points of $H(q)$ in $\Pi$. A point is \emph{free} if it is not in $\Pi$ and is at distance 5 from $v$; by
Lemma~\ref{lem:Piz}(d) there are $(q+1)(q^4+q^2+1)-(q+1)-q^2(q+1)-q^3(q+1)=q^3(q^2-1)$ of them.

\begin{definition}\label{def:whiskers}
A \emph{whisker set} is a set $W$ of lines of $H(q)$, each opposite $v$ and not in $S$, such that no free point is on two
lines of $W$. It is \emph{exact} if every free point is on exactly one line of $W$.
\end{definition}

A line opposite $v$ and not in $S$ meets $\Pi$ in one point, not on $v$ (Lemma~\ref{lem:Piz}(c),(d)), and its other $q$
points are free (every point of a line opposite $v$ is at distance 5 from $v$). So $|W|\le q^3(q^2-1)/q=q^2(q^2-1)$, with
equality if and only if $W$ is exact.

\begin{theorem}[doubling with whiskers]\label{thm:whiskerdoubling}
Let $W$ be a whisker set and $m\ge2$. Put
\[
\cA_2=\{v\},\qquad \cA_1=\{v\}\cup\cP(v)\cup S,\qquad \cA_0=\cA_1\cup\cP_\Pi\cup\cL_v\cup W .
\]
Then $\cA$ satisfies (L0), (L1) and (C$_6$). Hence $H(q)^{(m)}_\cA$ is bipartite with maximum degree $q+m$, diameter 6,
and
\[
2(q+1)(q^4+q^2+1)+(m-1)\bigl(q^4+2q^3+q^2+2q+2+|W|+m(q^3+q+2)+m^2\bigr)
\]
elements.
\end{theorem}

\begin{proof}
(L1), $i=1$: the elements on $v$ are its points (in $\cA_1$); a point of $v$ is on one element of $\cA_2$, a line of
$S\setminus\{v\}$ on none. (L1), $i=0$: the elements incident with $\cA_1$ are the points of $v$, the lines of $\cL_v$
and the points of $\Pi$, all in $\cA_0$; a line of $\cL_v$ meets $v$ in one point and contains no other point of $\cA_1$;
a point of $\cP_\Pi\setminus\cP(v)$ is on exactly one line of $S$ (Lemma~\ref{lem:Piz}(b)) and not on $v$; a whisker has
no point on $v$. (L0): a point outside $\cA_0$ is off $\Pi$. If it is at distance 3 from $v$, it is on exactly one line of
$\cL_v$ (two would close a circuit of length at most 6) and on no whisker (a whisker through it would be at distance at
most 4 from $v$). If it is free, it is on no line of $\cL_v$ and on at most one whisker. A line outside $\cA_0$ is not in
$\Pi$ (the lines in $\Pi$ are those of $S$), so it meets $\Pi$ in one point, its only point in $\cA_0$. (C$_6$): a line of
$\cL_v$ has no point of $\Pi$ other than its point on $v$ (Lemma~\ref{lem:Piz}(d)), and a whisker has exactly one point of
$\Pi$, not on $v$. So the components of the substructure on $\cA_0$ are
\[
\{v\}\cup\cP(v)\cup\cL_v\qquad\text{and}\qquad \{n\}\cup\cP(n)\cup\{h\in W:h\text{ meets }n\}\quad(n\in S\setminus\{v\}),
\]
trees of diameter 4 (the ball of radius 2 about $v$; a line, its points and the whiskers hanging from them).
Lemma~\ref{lem:hat} gives (C$_6$), and Theorem~\ref{thm:delorme} gives degree at most $q+1+(m-1)$, diameter at most 6 and
the count, with $|\cA_2|=1$, $|\cA_1|=(q^3+1)+(q+1)$ and $|\cA_0|=(q^3+1)+(q^2+q)+|W|+(q+1)(q^3+1)$. The diameter is
exactly 6: a free point $p$ (level $-1$) has $q^5$ opposite points, and only $(q+1)(q^3+1)<q^5$ points lie in $\cA_0$, so
some point $p'$ opposite $p$ has level $-1$, and $d(p,p')\ge6$ by projection.
\end{proof}

\emph{The neighbourhood of every version.} Write $x^w$ for the versions ($i,j,k\in[m]$). Every element of $H(q)$
appears once in Table~\ref{tab:versions}. An element of level $k$ has $m^{k+1}$ versions.

\begin{longtable}{@{}p{3.6cm} l p{6.2cm} l@{}}
\caption{The neighbourhood of every version in Theorem~\ref{thm:whiskerdoubling}.}\label{tab:versions}\\
\toprule element (level) & versions & neighbours of the version & degree\\ \midrule\endfirsthead
\toprule element (level) & versions & neighbours of the version & degree\\ \midrule\endhead
\bottomrule\endfoot
$v$ (2) & $v^{ijk}$ & $p^{ij}$ for the $q+1$ points $p$ of $v$ & $q+1$\\
point $p$ of $v$ (1) & $p^{ij}$ & $v^{ij1},\dots,v^{ijm}$; $L^i$ for the $q$ lines $L\in\cL_v$ on $p$ & $q+m$\\
line $n\in S\setminus\{v\}$ (1) & $n^{ij}$ & $p^i$ for the $q+1$ points $p$ of $n$ & $q+1$\\
point $p\in\cP_\Pi\setminus\cP(v)$ (0) & $p^i$ & $n^{i1},\dots,n^{im}$ for the line $n\in S$ on $p$; $h^i$ for each whisker $h$ on $p$; every other line on $p$ & $q+m$\\
line $L\in\cL_v$ (0) & $L^i$ & $p^{i1},\dots,p^{im}$ for $p=L\cap v$; the $q$ other points of $L$ & $q+m$\\
whisker $h\in W$ (0) & $h^i$ & $p^i$ for $p=h\cap\Pi$; the $q$ free points of $h$ & $q+1$\\
point $r\notin\Pi$ ($-1$) & $r$ & all $m$ versions of the line of $\cL_v\cup W$ on $r$ (if any); every other line on $r$ & $q+m$, or $q+1$ if $r$ is free and on no whisker\\
line $N\notin\cA_0$ ($-1$) & $N$ & all $m$ versions of its point in $\Pi$ (not on $v$); its $q$ other points & $q+m$\\
\end{longtable}

\subsection{An exact whisker set for every odd $q$}
For a point $X$ of $H(q)$ in $\Pi$, not on $v$, put
\[
\nu(X)=\frac{X_2X_3-X_4X_5}{cX_2^2+X_5^2}.
\]
The denominator vanishes only if $X_2=X_5=0$, that is (Lemma~\ref{lem:Piz}(d)) only on $v$. Let
\[
G(a,b)=(c_0a+c_1b)(a^2+cb^2)-b^3,\qquad F(z)=\frac{G(\alpha,\beta)}{\alpha^2+c\beta^2}\quad(z=(\alpha,\beta)\ne0),
\]
with $c_0,c_1\in\F$ chosen so that $G$ has no zero on $\F^2\setminus\{0\}$.

\begin{lemma}\label{lem:cubic}
Such $c_0,c_1$ exist for every odd $q$.
\end{lemma}

\begin{proof}
$G(a,0)=c_0a^3$, so $c_0\ne0$ is needed. For $b\ne0$ and $t=a/b$, $G(a,b)=b^3\bigl((c_0t+c_1)(t^2+c)-1\bigr)$, and
$t^2+c\ne0$. So $G$ has a zero with $b\ne0$ if and only if $(c_0,c_1)$ lies on one of the $q$ lines $c_0t+c_1=1/(t^2+c)$,
$t\in\F$, of the $(c_0,c_1)$-plane; each has $q-1$ points with $c_0\ne0$. The lines for $t=0$ and $t=1$ meet at
$c_0=1/(1+c)-1/c=-1/(c(1+c))\ne0$ ($c\ne-1$, as $-c$ is a non-square). So these lines cover at most $q(q-1)-1$ of the
$q(q-1)$ points with $c_0\ne0$.
\end{proof}

One valid choice for the $q$ of the table ($\F_9=\F_3[i]$, $i^2=-1$):
\begin{center}
\begin{tabular}{@{}l cccccc@{}}
\toprule $q$ & 3 & 5 & 7 & 9 & 11 & 13\\ \midrule
$c$ & 1 & 2 & 1 & $1+i$ & 1 & 2\\
$(c_0,c_1)$ & (1,2) & (1,0) & (1,0) & $(i,2i)$ & (1,3) & (1,1)\\ \bottomrule
\end{tabular}
\end{center}

\begin{theorem}[exact whiskers]\label{thm:whiskers}
Let $q$ be odd, $-c$ a non-square, and $c_0,c_1$ as in Lemma~\ref{lem:cubic}. Then
\[
W=\bigl\{\ell\in S_z\setminus\{v\}:\ z\in\F^2\setminus\{0\},\ \nu(\ell\cap\Pi)=F(z)\bigr\}
\]
is an exact whisker set: every free point lies on exactly one line of $W$, and $|W|=q^2(q^2-1)$.
\end{theorem}

Here $\ell\cap\Pi$ is one point not on $v$: $\ell\subseteq\Pi_z$ and $\ell\ne v$, so $\ell\notin S$ by
Lemma~\ref{lem:Piz}(e), and $\ell$ is opposite $v$ by Lemma~\ref{lem:Piz}(b). The proof uses a complex structure on each
$\Pi_z$. Write a point of $\Pi_z$ as $x=(x_1,\dots,x_6)$, meaning
$(X_0,\dots,X_6)=(-cx_4-\alpha x_2-\beta x_5,x_1,\dots,x_6)$; on $\Pi_z$ the quadratic form is
$Q_z(x)=x_1x_5+x_2x_6-x_3^2-cx_4^2-\alpha x_2x_4-\beta x_4x_5$, with polar form $b_z$. Put
$\mathrm{num}(y)=y_2y_3-y_4y_5$ and $\mathrm{den}(y)=cy_2^2+y_5^2$, and define the linear map
\[
J_zx=\bigl(-\beta x_3+\alpha x_4-x_6,\ x_5,\ -cx_4-\beta x_5,\ x_3+\beta x_2,\ -cx_2,\
cx_1+\alpha\beta x_2+\alpha x_3+c\beta x_4+\beta^2x_5\bigr);
\]
for $z=0$, $J_0x=(-x_6,x_5,-cx_4,x_3,-cx_2,cx_1)$.

\begin{lemma}[complex structure]\label{lem:Jz}
For every $z$:
\begin{enumerate}[label=(\alph*)]
\item $J_z^2=-c$, $Q_z(J_zx)=cQ_z(x)$ and $b_z(x,J_zx)=0$.
\item The six differences $p_{12}-p_{34}$, $p_{54}-p_{32}$, $p_{20}-p_{35}$, $p_{65}-p_{30}$, $p_{01}-p_{36}$, $p_{46}-p_{31}$ of
  the Grassmann coordinates of $x\wedge J_zx$ equal $Q_z(x)$ times $1,0,0,-c,\alpha,-\beta$.
\item For every point $x$ of $Q(6,q)\cap\Pi_z$, the line $\langle x,J_zx\rangle$ is the line of $S_z$ through $x$.
\item $\mathrm{num}(sx+tJ_0x)=(s^2+ct^2)\,\mathrm{num}(x)$ and $\mathrm{den}(sx+tJ_0x)=(s^2+ct^2)\,\mathrm{den}(x)$. So $\nu$
  is constant on every line of $S\setminus\{v\}$.
\end{enumerate}
\end{lemma}

\begin{proof}
(a), (b) and (d) are polynomial identities with integer coefficients, verified by direct expansion. For instance, for
$z=0$: $Q_0(J_0x)=cx_2x_6+cx_1x_5-c^2x_4^2-cx_3^2=cQ_0(x)$; $p_{12}=x_1x_5+x_2x_6$ and $p_{34}=x_3^2+cx_4^2$, so
$p_{12}-p_{34}=Q_0(x)$; in (d), $\mathrm{num}(J_0x)=c\,\mathrm{num}(x)$ and $\mathrm{den}(J_0x)=c\,\mathrm{den}(x)$, and the
polar forms $x_2y_3+x_3y_2-x_4y_5-x_5y_4$ and $2cx_2y_2+2x_5y_5$ vanish at $y=J_0x$. (c) $J_z$ has no eigenvector ($-c$ is
a non-square), so $x$ and $J_zx$ span a line; if $Q_z(x)=0$, then by (a) the whole line is singular and by (b) its
Grassmann coordinates satisfy \eqref{eq:tits}, so it is a line of $H(q)$ in $\Pi_z$ through $x$, the one of
Lemma~\ref{lem:Piz}(b). (d) By (c), a line $n\in S\setminus\{v\}$ is $\{sx+tJ_0x\}$ for any of its points $x$, and
$s^2+ct^2\ne0$.
\end{proof}

Let $p$ be a free point. By Lemma~\ref{lem:Piz}(d), $(p_2,p_5)\ne(0,0)$; as $p\notin\Pi$ we normalise $p_0+cp_4=1$. Then
$p\in\Pi_z$ if and only if $z$ lies on the affine line
\[
A_p=\{(\alpha,\beta):\ \alpha p_2+\beta p_5=-1\}\subseteq\F^2\setminus\{0\},
\]
which has $q$ points and direction $(p_5,-p_2)$. For $z\in A_p$ let $\ell_p(z)$ be the line of $S_z$ on $p$ and
$X(z)=\ell_p(z)\cap\Pi$.

\begin{lemma}[key identity]\label{lem:key}
For $z=(\alpha,\beta)\in A_p$,
\[
\Phi_p(z):=\nu(X(z))+\frac{\beta^3}{\alpha^2+c\beta^2}=\frac{p_2p_3-p_4p_5+\beta p_2^2}{cp_2^2+p_5^2}.
\]
So $\Phi_p$ is affine on $A_p$, with increment $-p_2^3/(cp_2^2+p_5^2)$ along $(p_5,-p_2)$.
\end{lemma}

\begin{proof}
Put $\Theta=cp_2^2+p_5^2\ne0$ and $\gamma=-\alpha p_5+c\beta p_2$. On $\Pi_z$ the linear form $\Lambda(y)=y_0+cy_4$ equals
$-\alpha y_2-\beta y_5$, and $\Pi=\{\Lambda=0\}$. We have $\Lambda(p)=1$ and $\Lambda(J_zp)=-\alpha p_5+c\beta p_2=\gamma$. By
Lemma~\ref{lem:Jz}(c), $\ell_p(z)=\langle p,J_zp\rangle$, so $X(z)=\gamma p-J_zp$, with
\[
X_2=\gamma p_2-p_5,\quad X_3=\gamma p_3+cp_4+\beta p_5,\quad X_4=\gamma p_4-p_3-\beta p_2,\quad X_5=\gamma p_5+cp_2 .
\]
Expanding, $\mathrm{den}(X)=(\gamma^2+c)\Theta$ and
$\mathrm{num}(X)=(\gamma^2+c)(p_2p_3-p_4p_5)+\beta(2\gamma p_2p_5+cp_2^2-p_5^2)$. Solving $\alpha p_2+\beta p_5=-1$,
$-\alpha p_5+c\beta p_2=\gamma$ (determinant $\Theta$) gives $\alpha=-(cp_2+\gamma p_5)/\Theta$ and
$\beta=(\gamma p_2-p_5)/\Theta$, hence $\alpha^2+c\beta^2=(\gamma^2+c)/\Theta\ne0$ and
$\beta^3/(\alpha^2+c\beta^2)=\beta(\gamma p_2-p_5)^2/((\gamma^2+c)\Theta)$. The two $\beta$-terms add up to
$\beta\bigl(2\gamma p_2p_5+cp_2^2-p_5^2+(\gamma p_2-p_5)^2\bigr)/((\gamma^2+c)\Theta)=\beta p_2^2/\Theta$. Along $(p_5,-p_2)$,
$\beta$ changes by $-p_2$.
\end{proof}

\begin{proof}[Proof of Theorem~\ref{thm:whiskers}]
Let $p$ be free. The lines of $\bigcup_{z\ne0}S_z\setminus\{v\}$ on $p$ are exactly the $q$ lines $\ell_p(z)$, $z\in
A_p$, and they are distinct (a common line would lie in $S_z\cap S_{z'}=\{v\}$); for the same reason each line of $W$
belongs to a single $S_z$. As $F(z)=c_0\alpha+c_1\beta-\beta^3/(\alpha^2+c\beta^2)$, we have $\ell_p(z)\in W$ if and only if
$\Phi_p(z)-(c_0\alpha+c_1\beta)=0$. By Lemma~\ref{lem:key} the left side is affine on $A_p$ with increment
\[
-\frac{p_2^3}{\Theta}-(c_0p_5-c_1p_2)=-\frac{G(p_5,-p_2)}{\Theta}\ne0
\]
along $(p_5,-p_2)$, so it vanishes at exactly one of the $q$ points of $A_p$: $p$ is on exactly one line of $W$. The lines
of $W$ are opposite $v$ and not in $S$. Counting incidences between free points and $W$ gives $|W|=q^3(q^2-1)/q$.
\end{proof}

\subsection{The bound}\label{ssec:thmD}

\begin{theorem}\label{thm:D}
For every odd prime power $q$ and every $m\ge2$, the structure of Theorem~\ref{thm:whiskerdoubling} with the exact whisker
set of Theorem~\ref{thm:whiskers} is bipartite, has maximum degree $q+m$ and diameter 6, and
\[
\nb(q+m,6)\ge2(q+1)(q^4+q^2+1)+(m-1)\bigl(2q^4+2q^3+2q+2+m(q^3+q+2)+m^2\bigr).
\]
In particular $\nb(q+2,6)\ge2(q+1)(q^4+q^2+1)+2q^4+4q^3+4q+10$ and
$\nb(q+3,6)\ge2(q+1)(q^4+q^2+1)+4q^4+10q^3+10q+34$.
\end{theorem}

\begin{proof}
Theorems~\ref{thm:whiskerdoubling} and~\ref{thm:whiskers}, with $|W|=q^4-q^2$.
\end{proof}

With $W$ exact, every element outside $\cA_0$ is incident with exactly one element of $\cA_0$. For a doubling of a
generalized $n$-gon of order $(q,q)$ with $N$ points, counting incidences gives $q|\cA_0|+2k\le2N$, $k$ the number of
components of $\cA_0$; here $k=q^3+1$ and $q(2q^4+2q^3+2q+2)+2(q^3+1)=2N$, so the single-level part is as large as this
counting bound allows for $q^3+1$ components.

\begin{remark}[Delorme's construction; the 2013 values]\label{rem:delorme6}
$W=\emptyset$ is Delorme's doubling of a generalized hexagon \cite[\S8(b)]{D85a}: he takes $q^3+1$ pairwise opposite
elements (citing \cite{T80}; here the Hermitian spread $S$), puts one of them in $\cA_2$, its neighbours and the others in
$\cA_1$, and their neighbours in $\cA_0$, so $|\cA_1|=q^3+q+2$ and $|\cA_0|=q^4+2q^3+q^2+2q+2$, exactly as in
Theorem~\ref{thm:whiskerdoubling} with $W=\emptyset$; at $q=5$ he states $7812+912+2\cdot132+4=8992$. At $m=2$ this has
$2(q+1)(q^4+q^2+1)+q^4+4q^3+q^2+4q+10$ elements: 8992, 142464, 374452, 842048 at $q=5,9,11,13$; at $q=13$, $m=3$, it has
884062. These are the 2013 values at $\Delta=7,11,13,15,16$. The exact whiskers add $(m-1)q^2(q^2-1)$.
\end{remark}

\emph{Worked counts.} $q=5$, $m=2$: $H(5)$ has $2\cdot6\cdot651=7812$ elements, $|\cA_0|=1512$ (of which $|W|=600$),
$|\cA_1|=132$, $|\cA_2|=1$, so the order is $7812+1512+2\cdot132+4=9592$. $q=13$, $m=3$: $H(13)$ has 804468 elements,
$|\cA_0|=61544$ ($|W|=28392$), $|\cA_1|=2212$, so the order is $804468+2(61544+3\cdot2212+9)=940846$.

\noindent\emph{Cells.} $m=2$: $(7,6)$ 9592 (2013: 8992), $(11,6)$ 148944 (142464), $(13,6)$ 388972 (374452), $(15,6)$ 870440
(842048). $m=3$, $q=13$: $(16,6)$ 940846 (884062). These graphs are not regular; see Section~\ref{ssec:thmR}.

\subsection{Even $q$}\label{ssec:evenq}
Let $q=2^e$, $\Tr$ the absolute trace of $\F$, and fix $\delta\in\F$ with $\Tr(\delta)=1$, so that $t^2+t+\delta$ is
irreducible and $a^2+ab+\delta b^2$ is anisotropic. Now $Q(6,q)$ is $X_0X_4+X_1X_5+X_2X_6+X_3^2=0$, with nucleus $e_3$,
and $H(q)$ is given by the same equations \eqref{eq:tits}. Put
\[
\Pi_z:\ X_3=X_0+\delta X_4+\alpha X_2+\beta X_5\qquad(z=(\alpha,\beta)\in\F^2),\qquad \Pi=\Pi_{(0,0)},
\]
$S_z$ the lines of $H(q)$ in $\Pi_z$, and $v=\langle e_1,e_6\rangle\subseteq\Pi_z$ as before.

\begin{lemma}\label{lem:Piz-even}
Lemma~\ref{lem:Piz} holds verbatim for even $q$ with these $\Pi_z$.
\end{lemma}

\begin{proof}
(a) On $\Pi_z$, $X_0X_4+X_1X_5+X_2X_6+X_3^2=(X_0^2+X_0X_4+\delta^2X_4^2)+(X_1+\beta^2X_5)X_5+X_2(X_6+\alpha^2X_2)$. The first
term is anisotropic ($\Tr(\delta^2)=\Tr(\delta)=1$) and the polar form is non-degenerate on $\Pi_z$ (which misses $e_3$), so
$\Pi_z\cap Q(6,q)$ is an elliptic quadric $Q^-(5,q)$. (b), (c), (e) are proved as in Lemma~\ref{lem:Piz}. (d) The same
three equations give $u_2=u_5=0$; on $X_2=X_5=0$ the quadric is the cone $X_0X_4=X_3^2$, and in $\Pi_z$ this becomes
$X_0^2+X_0X_4+\delta^2X_4^2=0$, so $X_0=X_4=X_3=0$, a point of $v$.
\end{proof}

So Theorem~\ref{thm:whiskerdoubling} holds for even $q$ without change. For $X\in\Pi$ not on $v$ put
\[
\nu(X)=\frac{X_2X_3+X_2X_4+X_4X_5}{\delta X_2^2+X_2X_5+X_5^2},\quad G(a,b)=(c_0a+c_1b)(a^2+ab+\delta b^2)+b^3,\quad
F(z)=\frac{G(\alpha,\beta)}{\alpha^2+\alpha\beta+\delta\beta^2}.
\]
The denominator of $\nu$ is anisotropic and vanishes only on $v$.

\begin{lemma}\label{lem:cubic-even}
There are $c_0,c_1\in\F$ with $G$ zero-free on $\F^2\setminus\{0\}$ if and only if $q\ge4$.
\end{lemma}

\begin{proof}
$G(a,0)=c_0a^3$ needs $c_0\ne0$. For $b\ne0$ and $t=a/b$: $G=0$ if and only if $c_0t+c_1=r(t):=1/(t^2+t+\delta)$, a line in
the $(c_0,c_1)$-plane with $q-1$ points having $c_0\ne0$. If $q\ge4$ take $\varpi\notin\{0,1\}$; then $r(\varpi)\ne r(0)$
(as $\varpi^2+\varpi\ne0$), so the lines for $t=0,\varpi$ meet at $c_0=(r(0)+r(\varpi))/\varpi\ne0$, and the $q$ lines cover
fewer than $q(q-1)$ points with $c_0\ne0$. If $q=2$ then $\delta=1$, $r(0)=r(1)=1$, and the lines $c_1=1$ and $c_0+c_1=1$
contain both points with $c_0=1$.
\end{proof}

On $\Pi_z$ use coordinates $x=(x_0,x_1,x_2,x_4,x_5,x_6)$, with $X_3=x_0+\delta x_4+\alpha x_2+\beta x_5$, and
\begin{multline*}
J_zx=\bigl((\delta+1)x_0+\delta\alpha x_2+\delta^2x_4+(\alpha+\delta\beta)x_5,\ \beta x_0+x_1+\alpha\beta x_2+(\alpha+\delta\beta)x_4+\beta^2x_5+x_6,\
x_2+x_5,\\ x_0+(\alpha+\beta)x_2+\delta x_4+\beta x_5,\ \delta x_2,\ (\alpha+\beta)x_0+\delta x_1+\alpha^2x_2+\delta\alpha x_4+\alpha\beta
x_5\bigr).
\end{multline*}

\begin{lemma}\label{lem:Jz-even}
With $Q$ the quadratic form and $b$ its polar form: (a) $J_z^2+J_z+\delta=0$, $Q(J_zx)=\delta Q(x)$ and $b(x,J_zx)=Q(x)$;
(b) the six differences of \eqref{eq:tits} for $x\wedge J_zx$ are $Q(x)$ times $1,0,0,\delta,\alpha,\beta$; (c) for every
point $x$ of $Q(6,q)\cap\Pi_z$, $\langle x,J_zx\rangle$ is the line of $S_z$ on $x$; (d) with
$\mathrm{num}=X_2X_3+X_2X_4+X_4X_5$ and $\mathrm{den}=\delta X_2^2+X_2X_5+X_5^2$ on $\Pi$,
$\mathrm{num}(sx+tJ_0x)=(s^2+st+\delta t^2)\mathrm{num}(x)$, and likewise for $\mathrm{den}$; so $\nu$ is constant on every
line of $S\setminus\{v\}$.
\end{lemma}

\begin{proof}
(a), (b), (d) are polynomial identities over $\F_2[\delta,\alpha,\beta,x]$, verified by expansion. (c) $J_z$ has no
eigenvector since $t^2+t+\delta$ is irreducible; if $Q(x)=0$ then $Q(sx+tJ_zx)=(s^2+st+\delta t^2)Q(x)=0$ by (a), and by (b)
the line satisfies \eqref{eq:tits}, so it is the line of Lemma~\ref{lem:Piz-even}(b). In (d), $s^2+st+\delta t^2\ne0$ for
$(s,t)\ne0$.
\end{proof}

Let $p$ be free; then $(p_2,p_5)\ne(0,0)$ by Lemma~\ref{lem:Piz-even}(d), and we normalise $p_0+p_3+\delta p_4=1$. With
$\Lambda(y)=y_0+y_3+\delta y_4$ we have $\Pi=\{\Lambda=0\}$ and $\Lambda=\alpha y_2+\beta y_5$ on $\Pi_z$, so $p\in\Pi_z$ if and
only if $z$ lies on the affine line $A_p:\ \alpha p_2+\beta p_5=1$ (direction $(p_5,p_2)$). For $z\in A_p$ let $\ell_p(z)$ be
the line of $S_z$ on $p$ and $X(z)=\ell_p(z)\cap\Pi$.

\begin{lemma}[key identity, characteristic 2]\label{lem:key-even}
For $z=(\alpha,\beta)\in A_p$,
\[
\Phi_p(z):=\nu(X(z))+\frac{\beta^3}{\alpha^2+\alpha\beta+\delta\beta^2}=\frac{p_2p_3+p_2p_4+p_4p_5+\beta
p_2^2}{\delta p_2^2+p_2p_5+p_5^2},
\]
which is affine on $A_p$ with increment $p_2^3/\Theta$ along $(p_5,p_2)$, $\Theta=\delta p_2^2+p_2p_5+p_5^2\ne0$.
\end{lemma}

\begin{proof}
$\Lambda(p)=1$ and $\gamma:=\Lambda(J_zp)=\alpha(p_2+p_5)+\delta\beta p_2$. By Lemma~\ref{lem:Jz-even}(c), $X(z)=\gamma p+J_zp$
(then $\Lambda(X)=2\gamma=0$); $X_2=(\gamma+1)p_2+p_5$ and $X_5=\delta p_2+\gamma p_5$. Expanding,
$\mathrm{den}(X)=(\gamma^2+\gamma+\delta)\Theta$ and
$\mathrm{num}(X)=(\gamma^2+\gamma+\delta)(p_2p_3+p_2p_4+p_4p_5)+\beta\bigl((\gamma+\delta+1)p_2^2+p_5^2\bigr)$, with
$\gamma^2+\gamma+\delta\ne0$. Solving $\alpha p_2+\beta p_5=1$, $\alpha(p_2+p_5)+\delta\beta p_2=\gamma$ (determinant $\Theta$):
$\alpha=(\delta p_2+\gamma p_5)/\Theta$ and $\beta=((\gamma+1)p_2+p_5)/\Theta$, so
$\alpha^2+\alpha\beta+\delta\beta^2=(\gamma^2+\gamma+\delta)/\Theta$ and
$\beta^3/(\alpha^2+\alpha\beta+\delta\beta^2)=\beta((\gamma+1)p_2+p_5)^2/((\gamma^2+\gamma+\delta)\Theta)$. The $\beta$-terms add to
$\beta\bigl((\gamma+\delta+1)+(\gamma+1)^2\bigr)p_2^2/((\gamma^2+\gamma+\delta)\Theta)=\beta p_2^2/\Theta$. Along $(p_5,p_2)$,
$\beta$ changes by $p_2$.
\end{proof}

\begin{theorem}[exact whiskers, even $q\ge4$]\label{thm:whiskers-even}
With $c_0,c_1$ as in Lemma~\ref{lem:cubic-even}, $W=\{\ell\in S_z\setminus\{v\}:z\ne0,\ \nu(\ell\cap\Pi)=F(z)\}$ is an
exact whisker set.
\end{theorem}

\begin{proof}
As for Theorem~\ref{thm:whiskers}: the lines of $\bigcup_{z\ne0}S_z\setminus\{v\}$ on a free point $p$ are the $q$
distinct lines $\ell_p(z)$, $z\in A_p$, and $\ell_p(z)\in W$ if and only if $\Phi_p(z)+c_0\alpha+c_1\beta=0$, since
$F(z)=c_0\alpha+c_1\beta+\beta^3/(\alpha^2+\alpha\beta+\delta\beta^2)$. The left side is affine on $A_p$ with increment
$p_2^3/\Theta+c_0p_5+c_1p_2=G(p_5,p_2)/\Theta\ne0$, so it vanishes at exactly one point.
\end{proof}

\begin{corollary}\label{cor:evenq}
Theorem~\ref{thm:D} and Theorem~\ref{thm:R} hold for every even $q\ge4$. For $q=2$ there is also an exact whisker set
(of $12=q^2(q^2-1)$ lines, given explicitly and computer-checked), so they hold for $q=2$ as well.
\end{corollary}

None of this changes a cell with $\Delta\le16$, for any $m$: for $q\in\{2,4,8\}$ and $\Delta=q+m\le16$, either $\Delta-1$ is
a prime power (a cell where the Moore bound is attained by $H(\Delta-1)$), or $\Delta\in\{7,11,13,15,16\}$, where
Theorem~\ref{thm:D} with the odd $q'=\Delta-2$ (or $q'=13$, $m=3$ at $\Delta=16$) gives much more (for instance 9592
against 4468 from $q=4$, $m=3$ at $\Delta=7$).
\todo{Corollary~\ref{cor:evenq}, $q=2$: the 12 whiskers are a computer-found set (its exactness is also proved in Lean).
List them in an appendix, or drop the $q=2$ case?}

\subsection{Regular graphs of diameter 6}\label{ssec:thmR}
The structure of Theorem~\ref{thm:D} has equally many points and lines in $\cA_0$ (exact whiskers force this), but $\cA_1$
has $q^3+1$ lines and only $q+1$ points, and $\cA_2$ is a line. So it has $(m-1)m(q^3-q+m)$ more lines than points. We
delete repeated lines and then add anti-flags.

\begin{lemma}[repeated lines]\label{lem:repeated}
In $H(q)^{(m)}_\cA$, for $n\in S\setminus\{v\}$ and $i\in[m]$, the $m$ lines $n^{i1},\dots,n^{im}$ all have the point set
$\{p^i:p\in n\}$; for $i,j\in[m]$ the lines $v^{ij1},\dots,v^{ijm}$ all have the point set $\{p^{ij}:p\in v\}$
(Table~\ref{tab:versions}). So there are $mq^3+m^2$ classes of $m$ repeated lines.
\end{lemma}

\begin{lemma}[completion]\label{lem:completion}
Let a finite incidence structure have maximum degree at most $\Delta$, every element of degree at least $\Delta-r$, and
$\sum_{\text{points}}(\Delta-\deg)=\sum_{\text{lines}}(\Delta-\deg)=:T$. If $T>2r\Delta$, there is a set $J$ of anti-flags
such that adding $J$ makes it $\Delta$-regular.
\end{lemma}

\begin{proof}
Choose a multiset $J$ of point--line pairs in which each element $x$ occurs $\Delta-\deg x$ times (possible as the sums
agree). Call an entry bad if it is a flag or a repeat of another entry, and take $J$ with the fewest bad entries. If
$(p,L)$ is bad: $L$ is joined (in $\I\cup J$, with multiplicity) to $\Delta$ points, each in at most $r$ entries of $J$;
likewise $p$. So at most $2r\Delta<T$ entries $(p',L')$ have $p'$ joined to $L$ or $L'$ joined to $p$, and some entry
$(p',L')$ has neither (it is not $(p,L)$ itself). Replacing $(p,L),(p',L')$ by $(p,L'),(p',L)$ removes a bad entry and
creates none. So the minimum is 0.
\end{proof}

\begin{theorem}\label{thm:R}
For every odd prime power $q$ and $m\ge2$ there is a $(q+m)$-regular incidence structure of diameter at most 6 with
\[
2(q+1)\bigl(q^4+q^2+1+(m-1)(q^3+m+1)\bigr)
\]
elements: $2(q+1)(q^4+q^3+q^2+4)$ for $m=2$ and $2(q+1)(q^4+2q^3+q^2+9)$ for $m=3$. For the five table instances below the
order exceeds $\Mb(q+m,5)$, so the diameter is exactly 6.
\end{theorem}

\begin{proof}
Start from Theorem~\ref{thm:D}. Fix a set $R$ of $q$ lines of $S\setminus\{v\}$, and delete the lines $v^{ijk}$ with
$k\ge2$ ($(m-1)m^2$ lines) and $n^{ij}$ with $j\ge2$ for $n\notin R$ ($(m-1)m(q^3-q)$ lines): in total
$(m-1)m(q^3-q+m)$, the excess, so points and lines are now equally many. Each class of Lemma~\ref{lem:repeated} keeps at
least one line, so by Lemma~\ref{lem:flags}(b) no distance between remaining elements increases. After the deletion the
degrees are: $p^{ij}$ ($p\in v$): $q+1$; $p^i$ on a line $n\notin R$ of $S$: $q+1$; $p^i$ on $n\in R$: $q+m$; $v^{ij1}$,
$n^{i1}$ ($n\notin R$), $n^{ij}$ ($n\in R$) and $h^i$ ($h\in W$): $q+1$; all others as in Table~\ref{tab:versions} ($q+m$).
So every deficit is at most $r=m-1$. Points and lines are equally many, so both deficit sums equal
$\Delta|\cP|-|E|=\Delta|\cL|-|E|$ ($\Delta=q+m$, $E$ the set of flags), as Lemma~\ref{lem:completion} requires; and their
common value $T$ is at least the deficit of the $m|W|$ whisker versions: $T\ge(m-1)mq^2(q^2-1)>2(m-1)(q+m)$.
Lemma~\ref{lem:completion} completes the structure; the added flags join a point to a line, so the graph stays bipartite,
and they do not increase distances. The number of points is $(q+1)(q^4+q^2+1)+(m-1)\bigl((q+1)(q^3+1)+m(q+1)\bigr)$.
\end{proof}

\emph{Neighbourhoods.} Every neighbourhood is that of Table~\ref{tab:versions}, minus the deleted lines listed in the
proof, plus the completing anti-flags: exactly $m-1$ of them at each element of degree $q+1$, none elsewhere. The
completing anti-flags are given by Lemma~\ref{lem:completion}, not by a formula; neither the maximum degree nor the
diameter depends on the choice.
\todo{Theorem~\ref{thm:R}: the completing anti-flags are only shown to exist (as in the Lean proof). If you want every
neighbourhood written out, we need an explicit completion (for instance a formula on the deficient elements); this is open
work.}

\noindent\emph{Cells (regular graphs).} $(7,6)$ 9348, $(11,6)$ 147500, $(13,6)$ 386328, $(15,6)$ 866068, $(16,6)$ 927724
($q=13$, $m=3$). Each beats the 2013 value of its cell (8992, 142464, 374452, 842048, 884062), and Delorme's graphs behind
those values are not regular (their point side is larger).

%%END%%

\bibliographystyle{amsplain}
\begin{thebibliography}{99}
%%BIB%%
\end{thebibliography}

\end{document}
